\documentclass[11pt]{article}

\usepackage[margin=1in]{geometry}
\usepackage{amsmath,amssymb,amsthm}
\usepackage{booktabs}
\usepackage{graphicx}
\usepackage{microtype}
\usepackage[colorlinks,linkcolor=blue,citecolor=blue]{hyperref}

\newtheorem{definition}{Definition}
\newtheorem{proposition}{Proposition}
\newtheorem{lemma}{Lemma}
\newtheorem{observation}{Observation}
\theoremstyle{remark}
\newtheorem{remark}{Remark}

\newcommand{\Yhat}{\widehat{Y}}
\newcommand{\Ylo}{Y_{\alpha}}
\newcommand{\Yhi}{Y_{1-\alpha}}
\newcommand{\HD}{\mathrm{HD}}
\newcommand{\HDF}{\mathrm{HD}_{95}}
\newcommand{\Dice}{\mathrm{Dice}}
\newcommand{\IoU}{\mathrm{IoU}}
\newcommand{\AURC}{\mathrm{AURC}}
\newcommand{\SDC}{\mathrm{SDC}}
\newcommand{\Ex}{\mathbb{E}}
\newcommand{\Prb}{\mathbb{P}}
\newcommand{\bd}{\partial}
\newcommand{\diag}{\mathrm{diag}}

\title{Metric-Aligned Confidence for Selective Segmentation:\\
Level-Set Posteriors and the Bi-Level Distance Band}
\author{}
\date{\today}

\begin{document}
\maketitle

\begin{abstract}
A selective segmentation system must decide, per image, whether to trust its own
mask or defer it to a human. This requires collapsing a per-pixel probability map
into one image-level confidence that ranks images by how bad the prediction
actually is. The right notion of ``bad'' is the deployment metric, and deployment
metrics disagree: overlap metrics (Dice, IoU) and boundary metrics (HD95, ASSD)
do not agree about which segmentations are failures. Existing post-hoc confidence
scores are almost all built for overlap risk; the strongest, Soft Dice Confidence
(SDC), is Dice-specific by construction.

We give one framework that produces a confidence aligned to \emph{any}
segmentation metric, post-hoc, from a single forward pass. We treat the
probability map as inducing a distribution over plausible ground truths --- the
\emph{level-set posterior}, the law of $\{i : p_i \ge T\}$ for
$T\sim\mathrm{Unif}(0,1)$, a classical random-set construction --- and score an
image by its expected loss under it. The expected loss collapses from a sum over
$2^N$ masks to a one-dimensional integral in the threshold, which we evaluate by
$M$-point quadrature. Instantiated at $M=2$ with $L=\HDF$, the estimator is a
single-pass, $O(N)$ score with a direct geometric reading: the width, in pixels,
of the probability band between the aggressive level set $\{p\ge\alpha\}$ and the
conservative one $\{p\ge 1-\alpha\}$, read by a Euclidean distance transform.

Our discriminating theoretical claim is a one-line lemma: SDC, and every
area-based readout of the band, is a \emph{permutation-invariant} functional of
$(p,\hat y)$ --- invariant under any joint relabelling of pixels --- and therefore
cannot resolve boundary geometry, whereas the distance readout can. This predicts,
before any experiment, that overlap risk should be well served by a spatially blind
score and boundary risk should not.

Across $4$ model conditions $\times\ 2$ datasets and $16$ published single-pass
baselines, the distance readout beats the band's own area readouts under boundary
risk in $8/8$ cells --- the paper's core result, and the only one invariant to every
convention we vary --- and beats SDC in $6/8$ under overlap risk and $6/8$ under a
boundary risk that charges detection failures rather than deleting them. Two further
results concern the estimator itself. At \emph{matched} nodes, the two-point quadrature
$r_2$ that the framework derives beats the band width we actually deployed on every
condition under every risk ($24/24$) --- same thresholds, same aggregation, the
functional the only difference --- at no extra cost; and denser quadrature beats SDC in
$7/8$ where the $M=2$ band manages $6/8$, contradicting our own $M=2$ argument. An audit
found an aggregation defect in the ladder that measured the latter and we withheld it;
it has since been recomputed, it reproduces to the digit, the defect never fired at the
deployed $\alpha$, and the embargo is lifted (Remark~\ref{rem:mladdercontam}).

We report the negative results rather than leaving them for a reviewer. The band's
win over SDC under boundary risk is $8/8$ only under the medical-imaging convention
that \emph{drops} images whose classes are present on one side; charging those
failures the image diagonal reduces it to $6/8$. The two cells it loses are the same
two it loses under overlap risk, which rules out a metric-alignment explanation: they
are CLIPSeg on multi-class VOC, where $54$--$67\%$ of images contain a hallucinated or
missed class, so the risk is dominated by \emph{detection} error --- and a
contour-displacement score cannot see whether a class is present at all. Two
\emph{spatially aware} baselines are also competitive under boundary risk, each
beating the band in one cell, which is itself evidence for the thesis. Finally, the
band's area readout is, up to reparameterization, the SAM stability score, which we
re-label as a published baseline rather than our own ablation.
\end{abstract}

\section{Introduction}

Segmentation models are deployed with a human in the loop. When the mask is good it
is accepted; when it is bad a clinician or annotator corrects it. The value of such
a system is governed not by average accuracy but by whether it can \emph{tell which
of its own outputs are bad}. This is selective prediction: the model emits, with its
mask, a scalar confidence $g(x)$; predictions are accepted in decreasing order of
$g$, and the system is judged by the risk--coverage trade-off.

We work in the \emph{post-hoc, image-level, single-pass} regime. Post-hoc: the
network is fixed and not retrained, because in practice one is handed a model.
Image-level: an entire image is accepted or rejected, because partial masks with a
few pixels abstained still require the expert to re-examine the whole image.
Single-pass: the score comes from one probability map --- no ensembles, no dropout
sampling, no test-time augmentation --- because these multiply inference cost.

The central question is how to collapse the $N$-dimensional probability map into
one number. The literature aggregates a per-pixel uncertainty --- mean maximum
softmax probability, mean entropy, patch and threshold variants. These share a
defect: they are stated in probability units and are agnostic to the metric the
application cares about. Soft Dice Confidence (SDC) \cite{sdc} fixed this for one
metric, showing the confidence should estimate the \emph{expected Dice of the
prediction under the label posterior}, and giving a near-optimal $O(N)$ estimator.

But Dice is not the only deployment metric, and often not the relevant one. A mask
can have excellent Dice and a catastrophic boundary --- a detached false-positive
blob far from the organ, a leak into an adjacent structure --- and boundary metrics
such as the $95$th-percentile Hausdorff distance (HD95) exist precisely to catch
what Dice does not. The two disagree strongly on our data: among VOC images with
\emph{identical} per-image mIoU, HD95 differs by up to $26$\,px
(\S\ref{sec:setup}), and the rank correlation between $1-\mathrm{mIoU}$ and HD95 is
only $0.46$--$0.67$ on VOC. A confidence aligned to Dice is, by construction,
aligned to the wrong thing when the risk is boundary error.

\paragraph{Contributions.}
\begin{enumerate}
\item \textbf{A metric-agnostic framework} (\S\ref{sec:framework}). We define the
ideal metric-aligned confidence $C^\star_L(x) = -\Ex[L(Y,\Yhat)\mid x]$ for an
arbitrary loss $L$, and make it computable by replacing the unknown label posterior
with the level-set posterior: the law of $Y^{(T)} = \{i : p_i \ge T\}$ under
$T\sim\mathrm{Unif}(0,1)$. Under it the ideal confidence collapses from $2^N$ masks
to a one-dimensional integral $\int_0^1 L(Y_t,\Yhat)\,dt$, evaluated by $M$-point
quadrature. The posterior itself is classical (\S\ref{sec:posterior}); what is new
is pairing it with a \emph{boundary} loss to obtain a post-hoc confidence.

\item \textbf{Theory: spatial blindness, not the coupling}
(\S\ref{sec:theory}). We prove SDC is the \emph{ratio-of-expectations} Dice
instantiation of our framework and is invariant to the coupling of the marginals
(Prop.~\ref{prop:sdc}). We then give the lemma that actually discriminates
(Lemma~\ref{lem:perm}): SDC and every area readout of the band are permutation
invariant in the pixel index, hence spatially blind, while the distance readout is
not. We are explicit that the stronger claim --- that expected Dice itself is
coupling-free --- is \emph{false}, and quantify it.

\item \textbf{The bi-level distance band} (\S\ref{sec:band}). The $M=2$, $L=\HDF$
instantiation: a Euclidean distance transform between the contours of
$\{p\ge\alpha\}$ and $\{p\ge 1-\alpha\}$, giving a confidence in \emph{pixels}, on
the scale of the metric it predicts, for two EDTs per present class.

\item \textbf{Experiments and five negative results} (\S\ref{sec:experiments}).
The distance-vs-area ablation is decisive and is the paper's core evidence. We also
report that the area readout coincides with the SAM stability score; that two
spatially aware baselines beat the band in one cell each; that $\alpha$ was
selected on the test set and the headline is not robust to it
(\S\ref{sec:threshold}); that vertex-aware node placement --- the rule our own
integrand analysis prescribes --- \emph{loses} to the nodes it was built to replace,
making it the fourth principled adaptive node rule in this paper to fail on measurement
(\S\ref{sec:openitems}); and that the two diagnostics we built to test our posterior's
assumption cannot identify the thing they were built to test, on this data or any
single-annotation data (\S\ref{sec:termb}).

\item \textbf{Threshold selection} (\S\ref{sec:threshold}). We show a fixed
threshold in probability space is structurally wrong --- for diffuse maps it
collapses the band onto the image diagonal outright on up to $17.7\%$ of images at
the deployed $\alpha=0.1$ and on $98.6$--$99.6\%$ at $\alpha=0.01$, and it pollutes
the multi-class score far more often than that: on $27.1$--$63.1\%$ of VOC images at
$\alpha=0.1$ at least one predicted class saturates, a rate fine-tuning does not fix
--- and give the corrected estimand and three label-free selection rules.
\end{enumerate}

\section{Setup}\label{sec:setup}

\paragraph{Selective segmentation.}
A fixed model maps an image $x$ to a probability map $p\in[0,1]^N$ over $N$ pixels
(one foreground class for now; \S\ref{sec:multiclass} handles $C$ classes) and to a
hard prediction $\Yhat = \{i : p_i \ge \gamma\}$, $\gamma=\tfrac12$. The ground
truth is a random mask $Y \sim \Prb(\cdot\mid x)$. A confidence $g$ and threshold
$\tau$ accept $x$ iff $g(x)\ge\tau$; for a loss $L\ge 0$,
\begin{equation}
\mathrm{cov}(\tau) = \Prb(g(x)\ge\tau),
\qquad
R(\tau) = \frac{\Ex[\,L(Y,\Yhat)\,\mathbf{1}\{g(x)\ge\tau\}\,]}{\Prb(g(x)\ge\tau)} .
\end{equation}
Sweeping $\tau$ traces the risk--coverage curve, summarized by
$\AURC(g)=\int_0^1 R(c)\,dc$, estimated as the mean selective risk over the $n$
empirical coverage levels $k/n$. Lower is better. One fact drives much of what
follows:
\begin{equation}\label{eq:rank}
\AURC(g) \text{ depends on } g \text{ only through the ranking it induces.}
\end{equation}
Hence any strictly increasing reparameterization of a score leaves AURC unchanged:
we never need the \emph{scale} of an estimate right, only its order.
(Ties are broken by a fixed, score-independent order, so \eqref{eq:rank} is a
statement about the induced total preorder; ties are not innocuous, and
\S\ref{sec:threshold} shows they are the mechanism by which a badly chosen
threshold destroys the score.)

\paragraph{Risks.} Three per-image risks, two of overlap type and one of boundary
type:
\begin{equation}
L_{\IoU} = 1 - \mathrm{mIoU}(Y,\Yhat), \qquad
L_{\Dice} = 1 - \mathrm{mDice}(Y,\Yhat), \qquad
L_{\HDF} = \HDF(Y,\Yhat),
\end{equation}
\begin{equation}\label{eq:hd95}
\HDF(A,B) = \mathrm{P}_{95}\Big(
\{ d(u,\bd B) : u\in\bd A \} \cup \{ d(v,\bd A) : v\in\bd B \}
\Big),
\quad d(u,S)=\min_{w\in S}\lVert u-w\rVert_2 ,
\end{equation}
with $\bd A$ the surface (mask minus its erosion). \textbf{Conventions, stated
because they matter.} Per-image mIoU/mDice are averaged over the classes observed
in the prediction or the ground truth of that image, \emph{background included}
(the VOC convention). HD95 is averaged over \emph{foreground} classes only, and is
undefined for a class present on one side only --- a hallucinated or a wholly
missed class. We therefore report two boundary conventions: $\HDF^{\mathrm{drop}}$
averages over classes present in both (the medical-imaging convention, which costs
a detection failure \emph{nothing} and deletes images that detect nothing:
$23$--$44\%$ of VOC class instances and up to $2.1\%$ of images), and
$\HDF^{\mathrm{pen}}$ charges each one-sided class the image diagonal, so every
risk is evaluated on the same image set. We report \emph{both} conventions throughout
(Tables~\ref{tab:boundary} and~\ref{tab:penalized}). The $\HDF^{\mathrm{pen}}$ run is
complete, and it \emph{degraded} the boundary headline: the band's win over SDC falls
from $8/8$ to $6/8$ (\S\ref{sec:results}, Finding~2), while the distance-vs-area
ablation is unchanged at $8/8$ under both. We report it degraded rather than reporting
only the convention that flatters us. (An earlier version of this paragraph described
$\HDF^{\mathrm{pen}}$ as ``pending re-run''; it is not, and \S\ref{sec:openitems}
closes the item.)

\paragraph{Overlap and boundary genuinely disagree.} On VOC we find images with
identical per-image mIoU and wildly different boundary error: at mIoU $0.333$,
HD95 of $0.00$ vs $26.03$\,px; at mIoU $0.885$, $7.10$ vs $21.31$\,px. Spearman
$\rho(1-\mathrm{mIoU},\HDF)$ is $0.46$--$0.67$ on VOC ($0.83$--$0.88$ on Pet). IoU
counts mislabelled pixels; it is blind to \emph{where} they are.

\section{Metric-aligned confidence}\label{sec:framework}

\subsection{The ideal confidence, and a bounded loss}

\begin{definition}[Ideal metric-aligned confidence]\label{def:ideal}
$C^\star_L(x) = -\,\Ex_{Y\sim\Prb(\cdot\mid x)}[\,L(Y,\Yhat(x))\,]$.
\end{definition}
With $L=1-\Dice$ this is the Ideal Dice Confidence of \cite{sdc}; the definition
removes the commitment to Dice. Ranking by $C^\star_L$ gives the pointwise-lowest
risk--coverage curve achievable by any function of $x$.

\begin{remark}[Bounded convention, required not optional]\label{rem:bounded}
$\HDF(\varnothing,\Yhat)$ is undefined, and the posterior below places strictly
positive mass on the empty mask. We therefore fix, once and for all,
\begin{equation}\label{eq:bounded}
L(\varnothing,\Yhat) := \diag \ \text{(boundary losses)}, \qquad
\Dice(\varnothing,\Yhat) := 0 ,
\end{equation}
$\diag$ the image diagonal. Without \eqref{eq:bounded} the expectation in
Definition~\ref{def:ideal} does not exist for $\HDF$. With it, $L$ is bounded and
every statement below is well posed. This convention is not cosmetic: it is the
sole reason $\Ex_{Q_p}[\HDF]$ is finite, and \S\ref{sec:threshold} shows it is
also what the estimator's failure mode is made of.
\end{remark}

Definition~\ref{def:ideal} is not computable: it is an expectation over $2^N$ masks
under a posterior a segmentation network never gives us. What the network gives is
the per-pixel marginals $p_i \approx \Prb(Y_i=1\mid x)$, and every post-hoc method
is a way of turning those into an estimate of Definition~\ref{def:ideal}.

\begin{remark}[We do not need the label; we need a law over labels]
The ideal confidence does not require the ground truth of the image at hand; it
requires a distribution over what it might be. A forecast of ``$70\%$ chance of
rain'' does not say whether it rains, yet suffices to compute the expected loss of
carrying an umbrella. The price is that the estimate inherits the quality of the
probability map: a miscalibrated model yields a miscalibrated confidence, and no
post-hoc score repairs it.
\end{remark}

\subsection{The level-set posterior}\label{sec:posterior}

The marginals do not determine the joint. We must \emph{choose} a joint consistent
with them, and we take the one generated by a single random global threshold.

\begin{definition}[Level-set posterior $Q_p$]\label{def:levelset}
Let $T\sim\mathrm{Unif}(0,1)$ and $Y^{(T)} = \{i : p_i \ge T\}$. Write $Q_p$ for its
law and $Y_t = \{i : p_i\ge t\}$ for the deterministic level set.
\end{definition}

\begin{remark}[This object is classical; we are importing it]\label{rem:classical}
$Q_p$ is the canonical \emph{coverage-function random set}: a random closed set
whose one-point coverage function is exactly $p$ \cite{goodman,molchanov}. Its two
defining properties --- (i) $\Prb_{Q_p}(Y_i=1)=p_i$, since $\Prb(T\le p_i)=p_i$;
and (ii) $\Prb_{Q_p}(Y_i=1,Y_j=1)=\min(p_i,p_j)$, the Fr\'echet--Hoeffding upper
bound, so $Q_p$ is the comonotone (inverse-transform, common-uniform) coupling of
the marginals, with almost surely nested samples --- are textbook. We claim no
novelty for them and cite the sources. The observation that an \emph{independent}
coupling destroys contour geometry, while a spatially coherent one preserves it, is
likewise established in the uncertain-isocontour literature
\cite{poethkow,bolin}. Our contribution begins one step later: pairing this
posterior with a \emph{boundary} loss to obtain a post-hoc, single-pass,
image-level confidence.
\end{remark}

Samples from $Q_p$ cost nothing --- they are thresholdings --- and are the nested
chain of level sets, $Y_1\subseteq\cdots\subseteq Y_0$. Low thresholds give
aggressive, dilated candidates; high thresholds conservative, eroded ones. Under
the independent coupling, a boundary region with $p\approx 0.5$ is realized as
per-pixel coin flips, i.e.\ salt-and-pepper speckle whose surface is fractal and
whose $\HDF$ against $\Yhat$ is dominated by isolated flipped pixels. Under $Q_p$
the same region is realized as a coherent inward or outward displacement of the
contour --- which is what segmentation ambiguity looks like.

\subsection{From posterior to score: a one-dimensional integral}

\begin{proposition}[Threshold-integral risk]\label{prop:integral}
For any loss $L$ bounded via \eqref{eq:bounded},
\begin{equation}\label{eq:integral}
r_{\mathrm{ls}}(x) := \Ex_{Y\sim Q_p}[L(Y,\Yhat)] = \int_0^1 L(Y_t,\Yhat)\,dt .
\end{equation}
\end{proposition}
\begin{proof}
$Y^{(T)}=Y_T$ with $T\sim\mathrm{Unif}(0,1)$, so the expectation over $Q_p$ is the
expectation over $T$. Boundedness (Remark~\ref{rem:bounded}) makes the integral
finite.
\end{proof}

Equation~\eqref{eq:integral} is exact \emph{for $Q_p$}, and holds for any $L$:
Dice, IoU, HD95, ASSD, surface Dice, or an application cost. This is the
metric-agnostic step. The integrand $H_x(t) = L(Y_t,\Yhat(x))$ is a scalar function
of one variable, and we estimate the integral by $M$-point quadrature:
\begin{equation}\label{eq:quad}
\boxed{\;
r_M(x) = \sum_{m=1}^{M} w_m\, L(Y_{t_m},\Yhat(x)),
\qquad
C_L(x) = -\, r_M(x)
\;}
\end{equation}
with $\sum_m w_m = 1$. Each term costs a thresholding plus one evaluation of $L$;
for $L=\HDF$ that is a Euclidean distance transform, $O(N)$.

\begin{remark}[The dense limit contains a max-softmax term]\label{rem:atom}
$Q_p$ places mass $1-\max_i p_i$ on the empty mask, so by \eqref{eq:bounded}
\begin{equation}
r_{\mathrm{ls}}(x) = \int_0^{\max_i p_i} H_x(t)\,dt \;+\; \big(1-\max_i p_i\big)\cdot\diag .
\end{equation}
The dense-$M$ reference therefore \emph{contains an explicit max-softmax term},
weighted by the image diagonal. On Pet, $\diag\approx 375$\,px against typical HD95
$\approx 12$\,px, so even $\max_i p_i = 0.97$ contributes $\approx 11$\,px. This is
a real property of the estimator, not an artifact: it is why we expect the dense
reference to disagree with $M=2$ precisely on the diffuse zero-shot conditions, and
it is the same mechanism that makes $\alpha=0.01$ degenerate
(\S\ref{sec:threshold}).
\end{remark}

\subsection{Where the error comes from}\label{sec:error}

Let $r_{\mathrm{true}}(x) = \Ex_{\Prb(\cdot\mid x)}[L(Y,\Yhat)]$. Then
\begin{equation}\label{eq:decomp}
| r_M - r_{\mathrm{true}} |
\;\le\;
\underbrace{| r_M - r_{\mathrm{ls}} |}_{\text{(a) quadrature error}}
\;+\;
\underbrace{| r_{\mathrm{ls}} - r_{\mathrm{true}} |}_{\text{(b) posterior model error}} .
\end{equation}
Term (a) is numerical and is controlled by $M$; \S\ref{sec:openitems} gives the
ablation that measures it. Term (b) is a modelling error and \emph{$M$ does not
reduce it}: it is the price of substituting $Q_p$ for the true posterior, and has
two sources --- whether $p$ is calibrated, and whether real ambiguity is nested
contour displacement rather than topological change. We do not claim (b) is zero.

We tried to measure it. \S\ref{sec:termb} reports that attempt; it is a negative result
about our own diagnostic, and it changes what the rest of the paper is allowed to say
about $Q_p$.

\subsection{Term (b) is not identified on this data --- a negative result about our own
diagnostic}\label{sec:termb}

We built two label-using diagnostics of term (b), and intended them as a headline:
\emph{our posterior's assumption nearly holds, SDC's is grossly violated everywhere.}
\begin{itemize}
\item \textbf{\texttt{levelset\_auc}} --- is the ground truth a level set of $p$? Rank
pixels by $p_i$ and ask how well that ranking separates $Y^*_i=1$ from $Y^*_i=0$. $Q_p$
predicts AUC $=1$ \emph{exactly}: under $Q_p$ every sample \emph{is} a level set of $p$.
\item \textbf{\texttt{residual\_moran\_i}} --- Moran's $I$ of the residual $Y^*-p$.
Conditional independence given $p$, which is SDC's assumption \cite[eq.~30]{sdc},
predicts $\approx 0$: independent residuals are spatially uncorrelated.
\end{itemize}
Measured across the eight conditions: \texttt{levelset\_auc} $\in[0.985, 0.999]$,
\texttt{residual\_moran\_i} $\in[0.841, 0.964]$. That reads exactly like the headline.

\textbf{Both statistics are implemented correctly, and the headline is not licensed by
them.} The failure is \emph{identification}, not power --- no sample size and no
sharpening of either statistic repairs it. We set it out in full because we believed the
headline for some time.

\paragraph{(i) Both statistics test a conjunction, and neither can attribute a
violation.} Write $q$ for the true marginal. The residual decomposes as
\begin{equation}\label{eq:residualdecomp}
Y^* - p \;=\; \underbrace{(Y^* - q)}_{\text{coupling}} \;+\; \underbrace{(q - p)}_{\text{calibration}} .
\end{equation}
A spatially structured $q-p$ --- a miscalibrated but spatially coherent map, which is
exactly what a segmentation network emits --- produces autocorrelated residuals
\emph{whatever the coupling is}. So Moran's $I$ tests the conjunction ``\emph{the
coupling is not independent} \textbf{and} $p=q$'', and a rejection cannot be assigned to
either leg. \texttt{levelset\_auc} tests the same conjunction from the other side.

\paragraph{(ii) A generator satisfying SDC's \emph{own} assumption reproduces the entire
measured signature.} We simulated the rival hypothesis: a deterministic mask $q$; ground
truths $Y^*_i \sim \mathrm{Bern}(q_i)$ drawn \emph{independently}, so that conditional
independence holds \textbf{exactly}, by construction; and $p = \mathrm{blur}(q)$ ---
accurate but miscalibrated. It yields \texttt{levelset\_auc} $=0.9995$ (including
$100\%$ of images at AUC exactly $1.0$, with a $6.3\%$ non-degenerate band),
\texttt{residual\_moran\_i} $\in[0.84, 0.95]$, $\rho(\text{auc},\text{mIoU})>0$ and
$\rho(\text{moran},\text{mIoU})<0$: every qualitative and quantitative feature we had
read as evidence \emph{for} $Q_p$ and \emph{against} independence.

\emph{Therefore ``Moran $0.84$--$0.96$ against $\approx 0$ under independence'' is not a
refutation of independence.} The $\approx 0$ null silently assumes perfect calibration,
which nobody has. We had computed our own side's null and never the rival's.

\paragraph{(iii) On Pet and VOC the question is vacuous, and the statistics are
measuring model accuracy.} A coupling is a statement about the joint law of an
\emph{ambiguous} label. Pet and VOC ground truth is near-deterministic --- there is no
real ambiguity about where the cat is --- and where the marginal is $\{0,1\}$-valued
\emph{every} coupling with those marginals coincides. There is nothing there for the
diagnostic to discriminate. What it does discriminate is how good the model is:

\begin{center}
\small
\begin{tabular}{lcc}
\toprule
statistic & $\rho$ with per-image mIoU, by condition & pooled \\
\midrule
\texttt{levelset\_auc}      & $+0.41$ to $+0.85$ & $\mathbf{+0.888}$ \\
\texttt{residual\_moran\_i} & $-0.28$ to $-0.77$ & $\mathbf{-0.49}$ \\
\bottomrule
\end{tabular}
\end{center}

An ``assumption check'' that correlates with mIoU at $0.888$ is an accuracy meter.

\paragraph{(iv) And we do not get to claim our own assumption holds either.} $Q_p$ makes
a point prediction --- \texttt{levelset\_auc} $=1.0$ exactly, with zero variance --- and
it is \textbf{refuted on $81\%$ of images}. The violation is tiny (median
$1-\mathrm{AUC} = 2\times 10^{-5}$), which is precisely why the aggregate reads
$0.985$--$0.999$ and looks like confirmation; but a point prediction with no variance is
refuted by any violation at all. The honest reading of our own table is that $Q_p$ is an
excellent approximation on this data \emph{and} false on it, and that we cannot say how
much of the residual is coupling error and how much is miscalibration.

\paragraph{The framing was worse than unsupported: it invented a tension that cannot
exist.} The story was ``SDC's assumption is violated, yet SDC survives''. But by
Prop.~\ref{prop:sdc} the SDC \emph{value} is a function of $(p,\hat y)$ alone --- no
coupling enters it --- so \emph{no coupling diagnostic can move it}. SDC is not
\emph{surviving} a violated assumption; it is \textbf{invariant} to it. (Its $<\!1\%$
error \emph{bound} is proved under independence and does not transport to $Q_p$ ---
Remark~\ref{rem:notcoupling} --- but that is a fact about the bound, not about the
score.) The subsidiary defence ``SDC is near-optimal under overlap risk anyway'' is also
false on our own data: it loses $6/8$ to the band under overlap risk, and its AURC is
$1.2$--$2.7\times$ the band's on all four Pet conditions against $0.8$--$1.1\times$ on
the four VOC ones (Table~\ref{tab:overlap}). It is near-optimal on VOC only.

\paragraph{What would decide it --- and why it cannot be run here.} The discriminating
data is \textbf{multiple annotations of the same image} (LIDC-IDRI, QUBIQ, multi-rater
lesion sets). Estimate $\hat q$ as the per-pixel mean over raters, then test the
\emph{spread of the foreground count}, on which the two couplings disagree by orders of
magnitude rather than in the third decimal:
\begin{equation}\label{eq:countspread}
\text{independence:}\quad \mathrm{Var}\Big(\sum_i Y_i\Big) = \sum_i q_i(1-q_i),
\qquad
\text{$Q_p$:}\quad \mathrm{Var}\Big(\sum_i Y_i\Big) \text{ at the Fr\'echet--Hoeffding maximum.}
\end{equation}
The independence null is closed-form and we validated it (Monte-Carlo sd $21.54$ against
analytic $21.42$); at matched sharpness $Q_p$ gives an sd $23$--$75\times$ larger ($830$
against $22$\,px). A handful of raters settles it, and it is exactly the quantity SDC's
own bound needs.

\textbf{Pet and VOC have a single ground truth per image, so this cannot be run on our
data at all.} That is the finding, and it is not a matter of effort: \emph{no reanalysis
of our evaluation can license either half of the headline.} It also changes the status
of the scope limitation in \S\ref{sec:limitations}. The absence of medical, multi-rater
data is not merely a gap in coverage --- it is a \textbf{precondition} for making any
claim about term (b) at all.

One cheaper interim check exists, and we have not run it: recalibrate $p$ (temperature,
or per-pixel isotonic on a held-out split) and recompute both statistics. If Moran's $I$
collapses toward $0$, the ``gross violation'' was miscalibration all along. Given
$\rho(\text{moran},\text{mIoU}) = -0.5$ to $-0.77$, we predict a large collapse.

\begin{remark}[The lesson, stated because we did not follow it]\label{rem:rivalnull}
A diagnostic that confirms your hypothesis has not confirmed it until you compute what
the \emph{rival} hypothesis predicts for the same statistic. Both couplings predicted
the same AUC and the same Moran's $I$. The entire finding was an artifact of only ever
computing our own side's null.
\end{remark}

\section{Relation to Soft Dice Confidence}\label{sec:theory}

SDC \cite{sdc} scores an image by the soft Dice between the probability map and the
hard prediction,
\begin{equation}\label{eq:sdc}
\SDC(p,\Yhat) = \frac{2\sum_i p_i \hat{y}_i}{\sum_i p_i + \sum_i \hat{y}_i},
\qquad \hat{y}_i = \mathbf{1}\{i\in\Yhat\},
\end{equation}
in $O(N)$, with no ground truth. It is the closest prior work.

\begin{proposition}[SDC is the ratio-of-expectations Dice instantiation]\label{prop:sdc}
Let $Q$ be \emph{any} coupling with marginals $p$. Then
\begin{equation}
\frac{\Ex_Q[\,2\sum_i Y_i\hat y_i\,]}{\Ex_Q[\,\sum_i Y_i + \sum_i \hat y_i\,]}
= \frac{2\sum_i p_i \hat y_i}{\sum_i p_i + \sum_i \hat y_i} = \SDC(p,\Yhat).
\end{equation}
Hence the ratio-of-expectations surrogate for expected Dice depends on the joint law
of $Y$ only through its marginals.
\end{proposition}
\begin{proof}
Numerator and denominator are linear in $Y$; linearity of expectation gives
$\Ex_Q[\sum_i Y_i\hat y_i]=\sum_i p_i\hat y_i$ and $\Ex_Q[\sum_i Y_i]=\sum_i p_i$
for any coupling with marginals $p$. The higher-order structure of $Q$ never
enters.
\end{proof}

\begin{remark}[What Prop.~\ref{prop:sdc} does \emph{not} say]\label{rem:notcoupling}
It is tempting --- and wrong --- to conclude that expected Dice is coupling-free.
$\Ex_Q[\Dice(Y,\Yhat)]$ is an expectation of a \emph{ratio}, a nonlinear function of
the two linear statistics, and the joint law of that pair does depend on the
coupling. Exhaustively, for $N=2$, $p=(\tfrac12,\tfrac12)$, $\Yhat=\{1\}$:
$\Ex_{\mathrm{indep}}[\Dice]=5/12$, $\Ex_{Q_p}[\Dice]=1/3$, $\SDC=1/2$ --- three
different values. Moreover the delta-method step cannot be rescued by
concentration: among all couplings with marginals $p$, the comonotone $Q_p$
\emph{maximizes} $\mathrm{Var}(\sum_i Y_i)$ (every covariance sits at its
Fr\'echet--Hoeffding maximum), so $Q_p$ is the coupling under which the Dice
denominator concentrates \emph{least}. SDC's own $<\!1\%$ relative-error bound is
proved under conditional independence \cite[eq.~30]{sdc} and does not transport to
$Q_p$. Under independence the bound is safe for the reason its authors give: $\sum_i
Y_i$ concentrates, so the expectation of the ratio approaches the ratio of expectations
at $O(1/N)$ --- a gap of $<\!0.4\%$ at the foreground masses in question. Under $Q_p$
that step never engages \emph{at any mass}: the variance sits at its
Fr\'echet--Hoeffding maximum by the sentence above, so the delta-method argument has
nothing to stand on and no amount of foreground mass rescues it. We state this plainly
because the discriminating argument does not need it.

\emph{A retraction.} Earlier versions of this remark quantified the $Q_p$ gap as
``$4.7\%$ mean / $26\%$ max, measured on synthetic maps''. \textbf{Those two numbers are
withdrawn.} No code in this project computes them; three downstream documents quoted
them back, two of the three having dropped the ``synthetic'' qualifier so that they read
as if they came from the real evaluation; and they are \emph{unidentifiable} in
principle --- across defensible synthetic families the mean spans $0.21$--$11.37\%$ and
the max $0.38$--$97.45\%$. The number was a property of an unrecorded generator, not of
$Q_p$, and no reconstruction of the generator could make it one. What remains above is a
\emph{derivation} ($<\!0.4\%$, from concentration) and a \emph{proof} (no shrinkage with
mass, from the Fr\'echet--Hoeffding maximum three lines earlier). Neither was ever a
measurement, and we no longer describe them as one.
\end{remark}

The property that actually does the work is spatial blindness, and it is immediate.

\begin{lemma}[Spatial blindness]\label{lem:perm}
Let $\sigma$ be any permutation of the pixel indices, acting jointly on $(p,\hat y)$.
Then $\SDC$, $\Ex_{Q_p}[\Dice]$, and the area readouts
$\Dice(\Yhi,\Ylo)$, $\IoU(\Yhi,\Ylo)$ of \S\ref{sec:area} are all invariant under
$\sigma$. The distance readout $\HDF(\Ylo,\Yhi)$ is not. Consequently no
permutation-invariant functional of $(p,\hat y)$ can resolve boundary geometry.
\end{lemma}
\begin{proof}
Each of $\sum_i p_i$, $\sum_i \hat y_i$, $\sum_i p_i\hat y_i$, $|\Ylo|$, $|\Yhi|$,
and the law of $|Y^{(T)}\cap\Yhat|$ under $Q_p$ is a sum over pixels and is
unchanged by relabelling. $\HDF$ depends on the Euclidean positions of the surface
pixels, which $\sigma$ does not preserve. A functional invariant under all
$\sigma$ is constant on the orbit of any $(p,\hat y)$, and that orbit contains both
a spatially coherent mask and a scattered one with the same boundary
\emph{cardinality} but arbitrarily different boundary \emph{geometry}.
\end{proof}

Verified numerically on a soft-boundary blob ($64\times64$, contour ramp
$p\!\in\![0.02,0.98]$), applying one joint permutation $\sigma$ to the pixel
indices of $(p,\hat y)$:

\begin{center}
\small
\begin{tabular}{lccc}
\toprule
score & original & permuted & invariant? \\
\midrule
$\SDC$                                   & $0.890158$ & $0.890158$ & yes \\
Band-Dice \eqref{eq:area}                & $0.713183$ & $0.713183$ & yes \\
SAM stability \eqref{eq:area}            & $0.554223$ & $0.554223$ & yes \\
Band-HD95 \eqref{eq:bandsym}             & $7.000$ & $1.414$ & \textbf{no} \\
$r_2$ \eqref{eq:r2}                      & $3.803$ & $1.000$ & \textbf{no} \\
\bottomrule
\end{tabular}
\end{center}

The three permutation-invariant scores are unchanged bit-for-bit; the two
distance scores collapse, because scattering the pixels destroys the coherent
contour whose displacement they measure --- while leaving every area and every
marginal statistic exactly intact.

\begin{remark}[What the lemma does and does not establish]\label{rem:permlimit}
An independent check confirms the mechanics exactly: under a joint permutation, $\SDC$,
$\Ex_{Q_p}[\Dice]$ and both area readouts are bit-for-bit unchanged while
$\HDF(\Ylo,\Yhi)$ moves ($7.616 \to 3.000$ on its construction), and the overlap
\emph{risk} $1-\mathrm{mIoU}$ is likewise unchanged ($0.18317$ both). That last fact is
not a sting but a confirmation: the overlap risk is blind to exactly what SDC is blind
to, which is why a spatially blind estimator suffices for it.

The \emph{limitation} is in the proof's reach. It exhibits two inputs in the same
permutation orbit --- one spatially coherent, one scattered --- with equal
$(\SDC, \text{areas})$ and different boundary geometry. But a scattered probability map
is not something a segmentation network emits: real inputs live on a thin manifold of
spatially coherent maps that the orbit essentially never visits. So the lemma is a
valid impossibility result \emph{over the full input space} $[0,1]^N\!\times\!\{0,1\}^N$,
and it is \emph{suggestive} rather than decisive for the realistic sub-manifold. What
makes the conclusion bite is the measurement, not the proof: on real maps SDC's
boundary AURC is uniformly poor (\S\ref{sec:results}). We state the lemma as the reason
to \emph{expect} that, not as a demonstration of it.
\end{remark}

\paragraph{The falsifiable prediction.} Lemma~\ref{lem:perm}, with its reach bounded as
in Remark~\ref{rem:permlimit}, predicts, before any
experiment: \emph{SDC should be near-optimal under overlap risk} --- an overlap
metric is itself a permutation-invariant functional of $Y$, so a spatially blind
estimator loses nothing essential --- \emph{and mediocre under boundary risk}, where
a spatially blind score is structurally unable to see the quantity being predicted.
Note what it does \emph{not} predict: it does not say SDC should \emph{beat} a
geometric score under overlap risk, only that it should not be beaten \emph{for
lack of geometry}. \S\ref{sec:results} confirms the prediction, including the
finer point that our own area readout --- a permutation-invariant score by
Lemma~\ref{lem:perm} --- tracks SDC.

\begin{remark}[What we inherit from \cite{sdc}]
The principle (confidence should estimate the expected deployment metric) and the
regime (post-hoc, single-pass, image-level). Not the guarantees: SDC's bounds follow
from the algebra of Dice and do not transfer to $\HDF$. Our guarantee is of a
different kind --- exactness under a stated posterior (Prop.~\ref{prop:integral})
plus the error decomposition \eqref{eq:decomp}.
\end{remark}

\section{The bi-level distance band}\label{sec:band}

\subsection{Two thresholds}

Fix $\alpha<\tfrac12$ and take $t_1=\alpha$, $t_2=1-\alpha$, $w_1=w_2=\tfrac12$.
The nodes bracket $\gamma=\tfrac12$ and give the nested triple
\begin{equation}\label{eq:nesting}
\Yhi = \{p \ge 1-\alpha\} \;\subseteq\; \Yhat = \{p\ge\tfrac12\} \;\subseteq\;
\Ylo = \{p\ge\alpha\} .
\end{equation}
$\Yhi$ is the conservative candidate, $\Ylo$ the aggressive one, and
$\Ylo\setminus\Yhi$ is the \emph{probability band} --- the region the model is
undecided about. The $M=2$ risk is
\begin{equation}\label{eq:r2}
r_2(x) = \tfrac12\Big[\, \HDF(\Ylo,\Yhat) + \HDF(\Yhi,\Yhat) \,\Big].
\end{equation}

\subsection{The band width, and its honest relation to $r_2$}

In practice we compute a cheaper quantity: the width of the band itself, bypassing
$\Yhat$,
\begin{equation}\label{eq:bandsym}
\widehat{\HDF}^{\,\mathrm{sym}}(x) = \HDF(\Ylo,\Yhi),
\end{equation}
two Euclidean distance transforms, $O(N)$. Directed variants read only one side:
$\widehat{\HDF}^{\,\mathrm{out}}$ (from $\bd\Ylo$ to $\bd\Yhi$; catches fog bulges
and detached mid-confidence blobs) and $\widehat{\HDF}^{\,\mathrm{in}}$ (the
reverse; blind to detached blobs). The confidence is any strictly decreasing map,
e.g.\ $C^{\mathrm{band}} = \exp(-\widehat{\HDF}^{\,\mathrm{sym}}/\sigma)$; by
\eqref{eq:rank}, $\sigma$ affects absolute thresholds only, never the ranking.

\begin{proposition}[Band width vs.\ two-point quadrature, for max-Hausdorff]\label{prop:band}
Let $\HD$ be the \emph{max} Hausdorff distance and
$r_2^{\HD} = \tfrac12[\HD(\bd\Ylo,\bd\Yhat)+\HD(\bd\Yhat,\bd\Yhi)]$. Then
\begin{equation}
\HD(\bd\Ylo,\bd\Yhi) \;\le\; 2\, r_2^{\HD}(x).
\end{equation}
\end{proposition}
\begin{proof}
$\HD$ is a metric on compact sets; apply the triangle inequality.
\end{proof}

\begin{remark}[Prop.~\ref{prop:band} motivates but does not establish the bridge]\label{rem:nobridge}
Three caveats, stated because each is a real limitation.
\begin{enumerate}
\item \textbf{Wrong functional.} Everything implemented and evaluated is $\HDF$, a
\emph{percentile} of pooled surface distances. $\HDF$ is \emph{not a metric} and
does not obey the triangle inequality, so Prop.~\ref{prop:band} does not transfer.
It can fail badly: on a $400\times400$ map (disk at $p=1$, thin spike at $p=0.7$,
second thin spike at $p=0.3$, $\alpha=0.1$, nesting verified) we get
$\HDF(\Ylo,\Yhat)=\HDF(\Yhi,\Yhat)=0$, so $2r_2^{\HDF}=0$, while the implemented
band width is $13.41$\,px --- violated by an unbounded factor. Percentile
truncation is not subadditive: each spike is under $5\%$ of the pooled surface in
its pairwise comparison and is truncated away, but together they survive.
\item \textbf{No equality condition.} We claim none. Even for max-$\HD$, equality
requires \emph{co-location} of the extremal displacements, which fails under any
asymmetric level-set growth; on $398$ realistic nested triples the equality
frequency was $0/398$ (mean ratio $\HD/2r_2 = 0.665$).
\item \textbf{No rank-equivalence.} A one-sided inequality that is tight on some
images and slack on others \emph{permutes} the ranking; it is not benign for a pure
rank statistic. On synthetic nested triples Kendall $\tau(\text{band}, r_2)=0.81$
with $9.5\%$ discordant pairs.
\end{enumerate}
The band width is therefore best described as a posterior-\emph{spread} proxy that
Prop.~\ref{prop:band} motivates, not as the two-point expected-loss estimate. We
have implemented $r_2$ directly; the comparison of their rankings is the ablation that
must settle this, it \emph{has} been run, and the bridge holds \emph{empirically} ---
$\rho(\text{band}, r_2)\in[0.974, 0.992]$ across all eight conditions
(\S\ref{sec:openitems}, item 2). That is a measurement, not a theorem, and it is the
only thing connecting the deployed band to the estimator the theory derives.
\end{remark}

\subsection{The area counterpart, and what it already is}\label{sec:area}

The same frame admits \emph{region} readouts. Because the sets are nested these
collapse to area ratios:
\begin{equation}\label{eq:area}
\IoU(\Yhi,\Ylo) = \frac{|\Yhi|}{|\Ylo|},
\qquad
\Dice(\Yhi,\Ylo) = \frac{2|\Yhi|}{|\Yhi|+|\Ylo|} .
\end{equation}

\begin{remark}[The area readout is the SAM stability score]\label{rem:sam}
$\IoU(\Yhi,\Ylo)$ is, up to reparameterization, the \emph{stability score} of
Segment Anything \cite{sam}, which computes the IoU between masks obtained by
thresholding the predicted logits at a high and a low value, notes in code that
``one mask is always contained inside the other,'' and uses the result at inference
as a per-mask confidence filter. With SAM's defaults the logit bracket maps through
the sigmoid to probability thresholds $0.7311/0.2689$, which sum to exactly $1$ ---
i.e.\ SAM's score \emph{is} $|\{p\ge 1-a\}|/|\{p\ge a\}|$, and our $\alpha$ is SAM's
\texttt{threshold\_offset} reparameterized. We therefore re-label it as a published
baseline rather than our own ablation. What remains ours is (i) the level-set
posterior / expected-loss \emph{interpretation} of that construction, and (ii)
replacing the area readout with a \emph{distance} readout aligned to boundary risk.
\end{remark}

\begin{remark}[The two area readouts are rank-identical in the binary case]
For a single foreground class the sets are nested, so
$\Dice = 2\,\IoU/(1+\IoU)$ is a strictly increasing transform of $\IoU$; by
\eqref{eq:rank} the two give \emph{identical} AURC. They can differ only through
the present-class mean in the multi-class case. This is why the Pet rows of
Tables~\ref{tab:overlap}--\ref{tab:boundary} are identical for the two, and it must
be disclosed rather than presented as two independent scores.
\end{remark}

Comparing \eqref{eq:bandsym} against \eqref{eq:area} holds the posterior, the
thresholds, and the aggregation fixed and varies \emph{only} whether the band is
read spatially or by area. By Lemma~\ref{lem:perm} the area form is permutation
invariant and the distance form is not, so this is a controlled ablation of exactly
the quantity the theory says should matter. It is the paper's core evidence. We also
report a cruder ablation, band area over predicted perimeter (``area/perim''), a mean
width using no distance transform.

\subsection{Multi-class aggregation, conventions, cost}\label{sec:multiclass}

We apply the score per foreground class and aggregate over the classes predicted
present, $\mathcal{P}(x)$, by the mean (a $\max$ variant is also reported).
Restricting to predicted-present classes matters: aggregating over all classes lets
an absent class with mid-range probabilities hijack the score.

\paragraph{Degenerate cases, and a bug they caused.} If $\Ylo=\varnothing$ the class
is confidently absent and contributes nothing. If $\Ylo\ne\varnothing$ but
$\Yhi=\varnothing$ the band has no inner anchor and the width saturates at $\diag$.
\textbf{If no foreground class is predicted at all}, the image is a \emph{total
detection failure} and must receive the \emph{lowest} confidence under every score.
Because the distance scores are negated, aggregating an empty present-class list to
$0$ would make these images the \emph{most} confident of the split. An earlier
version of our code did exactly this; the images affected have mean $1-\mathrm{mIoU}$
of $0.55$--$0.84$ against dataset means of $0.04$--$0.40$, and the resulting AURC
error is large enough to invert conclusions (\S\ref{sec:results}). We now assign
them $-\diag$, matching SDC's convention (whose $0$ is its minimum). We flag this
because it is exactly the kind of convention that silently decides a headline.

\paragraph{Cost.} One forward pass, then $2C$ thresholdings and $2C$ Euclidean
distance transforms: $O(CN)$, milliseconds per image, no retraining, no auxiliary
model. Note $\alpha$ \emph{does} require a calibration set (\S\ref{sec:threshold});
we do not claim otherwise.

\section{Thresholds: the estimator's weakest joint}\label{sec:threshold}

A fixed $\alpha$ in probability space carries a structural risk. The condition for
a usable band is $\Yhi\ne\varnothing$, i.e.\ $1-\alpha \le \max_i p_i$ --- a
constraint on each \emph{image's own upper tail}, which a global constant has no
reason to satisfy. Diffuse (zero-shot, sigmoid) maps rarely reach $0.9$, so their
band saturates at $\diag$.

\paragraph{How bad it actually is, and on which counting rule.} Degeneracy admits
two rates, and the difference between them is the whole argument below, so we
report both. \emph{Total} degeneracy (columns $T$) is the fraction of images whose
band \emph{is} the image diagonal: every class the model predicts present has an
empty conservative set, so the score carries no ranking information at all.
\emph{Partial} degeneracy (columns $P$) is the fraction of images on which
\emph{at least one} present class saturates: the present-class mean of
\S\ref{sec:multiclass} is then polluted by a $-\diag$ term, and the score is
corrupted rather than destroyed. The two coincide whenever a single class is
predicted --- which is why they are \emph{identical} on all four binary Pet
conditions and diverge only on VOC. Both are measured against each image's true
native-resolution diagonal, not a proxy, and both count only saturation: the
$0.0$--$1.1\%$ of images that floor at $-\diag$ because they detect \emph{nothing}
(\S\ref{sec:multiclass}) are excluded from both, being a different failure.

\begin{center}
\small
\begin{tabular}{lcccccc}
\toprule
& \multicolumn{2}{c}{$\alpha=0.01$} & \multicolumn{2}{c}{$\alpha=0.05$}
& \multicolumn{2}{c}{$\alpha=0.1$} \\
\cmidrule(lr){2-3}\cmidrule(lr){4-5}\cmidrule(lr){6-7}
condition & $T$ & $P$ & $T$ & $P$ & $T$ & $P$ \\
\midrule
clipseg-general, Pet    & \textbf{99.6\%} & 99.6\% & 37.9\% & 37.9\% & 14.7\% & 14.7\% \\
clipseg-general, VOC    & \textbf{98.6\%} & 99.3\% & 46.3\% & 79.8\% & 17.7\% & \textbf{63.1\%} \\
clipseg-target, Pet     & 0.3\% & 0.3\% & 0.1\% & 0.1\% & 0.0\% & 0.0\% \\
clipseg-target, VOC     & 17.3\% & 63.6\% & 5.1\% & 47.2\% & 1.4\% & \textbf{35.9\%} \\
deeplabv3-external, Pet & 4.6\% & 4.6\% & 1.8\% & 1.8\% & 1.1\% & 1.1\% \\
deeplabv3-external, VOC & 18.4\% & 47.7\% & 8.4\% & 37.1\% & 5.0\% & \textbf{32.1\%} \\
deeplabv3-target, Pet   & 0.1\% & 0.1\% & 0.0\% & 0.0\% & 0.0\% & 0.0\% \\
deeplabv3-target, VOC   & 12.0\% & 41.3\% & 5.4\% & 31.5\% & 2.7\% & \textbf{27.1\%} \\
\bottomrule
\end{tabular}
\end{center}

\emph{A note for anyone checking these against the code}, because we got it wrong
ourselves: neither column is the mean of the \texttt{diag\_saturated@} diagnostic
that \texttt{scripts/show\_results.py} prints. That diagnostic is a per-image
\emph{share of present classes}, and its mean over images ($41.8\%$ on
\texttt{clipseg-general}/VOC at $\alpha=0.1$) is a third quantity, bounded between
$T$ and $P$ and equal to both only in the binary case. $T$ is
$\texttt{diag\_saturated@}=1$ with a present class; $P$ is
$\texttt{diag\_saturated@}>0$.

At $\alpha=0.01$ the zero-shot CLIPSeg score is dead: $98.6$--$99.6\%$ of images
saturate outright, and the score is $92\%$ tied ($3669$ images take only $511$
distinct values). This, and not a smooth loss of sensitivity, is what the
$\alpha=0.01$ column of any ablation is measuring.

At the deployed $\alpha=0.1$ the two rates tell different stories, and they are
governed by \emph{different} factors, which is why reporting only one is
misleading. \emph{Total} degeneracy is confined to $0$--$17.7\%$ of images, and it
tracks \textbf{how diffuse the map is}, exactly as the mechanism above predicts:
$14.7\%$ and $17.7\%$ on zero-shot CLIPSeg, whose sigmoid maps rarely reach $0.9$,
against $0.0$--$5.0\%$ on every other condition and $0.0$--$2.7\%$ once the model is
fine-tuned. On $T$, then, the degeneracy really is essentially absent on the
fine-tuned conditions. \emph{Partial}
degeneracy tracks something else entirely --- \textbf{the number of classes} --- and
fine-tuning does not touch it. It is $0.0$--$14.7\%$ on Pet, where it coincides with
$T$ exactly because there is only one foreground class to saturate; but it is
$27.1$--$63.1\%$ on \emph{all four} VOC conditions, fine-tuned ones included
($35.9\%$ on \texttt{clipseg-target}/VOC, $27.1\%$ on \texttt{deeplabv3-target}/VOC,
against total rates of $1.4\%$ and $2.7\%$). With $C=21$ the present-class mean
almost always has some weak class in it, however well trained the model is.

\emph{An earlier version of this paragraph read: ``the degeneracy is confined to
$0$--$17.7\%$ of images, and it is essentially absent on the fine-tuned
conditions.''} Both clauses are true --- of $T$. Neither is true of $P$, which we had
not computed. The sentence was therefore correct as written and wrong as read, and
it under-reported the failure by a factor of up to $26$ on the cell that matters most
(\texttt{clipseg-target}/VOC: $1.4\%$ against $35.9\%$). We flag it because the
smaller number is the one that supported our own argument, and Remark~\ref{rem:retraction}
is where it did the damage.

\begin{remark}[A retraction, and a partial retraction of the retraction]\label{rem:retraction}
An earlier draft of this section reported saturation rates of $27$--$67\%$ at
$\alpha=0.1$ and argued that our reported results therefore \emph{understate} the
method, with the largest gains to come on VOC. The rates were measured with a
width-$>200$\,px proxy rather than against each image's own diagonal, and we
retracted them as an over-count. \textbf{That retraction was half right, and we
correct it here rather than leave the convenient version standing.}

What was right: the proxy is not a measurement of $T$. Against each image's true
diagonal, total degeneracy at $\alpha=0.1$ is $1.4$--$17.7\%$ on VOC, not
$27$--$67\%$, and the claim that the method's headline is being suppressed by
wholesale band collapse does not survive.

What was wrong: we described the proxy as over-counting and moved on, when in fact
the $27$--$67\%$ band is a close match to a real quantity we had simply not
computed --- the \emph{partial} rate $P$, which is $27.1$--$63.1\%$ across the four
VOC conditions. A band wider than $200$\,px is not evidence that the band \emph{is}
the diagonal, but it is decent evidence that a $-\diag$ term has entered the mean.
The proxy was measuring the right phenomenon with the wrong estimator, and we
retracted the phenomenon along with the estimator.

The inference the retraction drew must therefore be redone, because it rested on
the smaller number. We had written: \texttt{clipseg-target}/VOC, one of the two
cells where the band \emph{loses} to SDC under overlap risk, is only $1.4\%$
saturated, \emph{so that loss cannot be a threshold artifact}. The premise is
correct and the inference does not follow --- at $P = 35.9\%$ that cell is not
obviously clean, and ``far too small'' is not an argument we can still make.

The conclusion survives, but it needs a different argument, and we must first concede
what the observational pattern cannot do. Ranked by $P$, the four VOC conditions are
$63.1$, $35.9$, $32.1$, $27.1\%$ --- and the two that lose to SDC are exactly the two
highest. Degeneracy and the nesting leak are \emph{collinear} across our conditions,
because CLIPSeg-on-VOC is worst at both. No cross-condition correlation can separate
them, and we do not claim one does.

What separates them is an \textbf{intervention}. The repair
\eqref{eq:mband} intersects the conservative set with the class's own prediction. It
targets the leak, and it \emph{provably cannot touch degeneracy}: if
$\{p_c \ge 1-\alpha\}$ is empty then so is its intersection with $\Yhat_c$, so a
saturated class stays saturated and its width stays $\diag$. We verify this rather
than assert it --- band and \eqref{eq:mband} are bit-identical on every saturated
image of all eight conditions. Yet the repair closes $42$--$77\%$ of the band's AURC
gap to SDC on precisely the two cells at issue (\S\ref{sec:results}). A
threshold-degeneracy explanation predicts that a leak-only intervention does
\emph{nothing}; it does most of the work. That is the discriminating evidence, and
it is the argument the retraction should have made --- not an appeal to a rate being
small.
\end{remark}

Saturated images do not, at $\alpha=0.1$, form one tied block: they collapse onto
$-\diag$, which varies with image size ($171$--$2596$\,px on Pet). Within the
block they are therefore ordered \emph{by image size}, and
$\rho(\diag, \text{true risk}) \approx 0.05$ --- i.e.\ ordered essentially at
random. That is the defect: not ties, but noise.

\paragraph{The estimand should be trimmed, and $\alpha$ is part of it.} By
Remark~\ref{rem:atom}, $H_x(t)$ saturates at $\diag$ as $t\to 1$ and blows up as
$t\to 0$ ($Y_t\to$ the whole image). The untrimmed
$\int_0^1 \HDF(Y_t,\Yhat)\,dt$ is dominated by degenerate level sets that nobody
would call plausible segmentations. The honest estimand is therefore the
\emph{trimmed} one,
\begin{equation}\label{eq:trimmed}
r_\alpha(x) = \Ex\big[\,L(Y_T,\Yhat)\;\big|\; T\in(\alpha,1-\alpha)\,\big],
\end{equation}
which makes $\alpha$ a \emph{definitional} significance level (``what counts as
plausible''), not a nuisance hyperparameter. The cost must be stated: the trimmed
posterior's marginals are the winsorized
$\mathrm{clip}\big((p_i-\alpha)/(1-2\alpha),\,0,\,1\big)$, so
Remark~\ref{rem:classical}(i) holds only in the trimmed sense.

\paragraph{The integrand's own order statistics are the exact nodes.} $H_x(t)$ is
\emph{piecewise constant} in $t$, jumping only where $t$ crosses one of the image's
own probability values. Hence, exactly,
\begin{equation}\label{eq:orderstat}
\int_0^1 H_x(t)\,dt \;=\; \sum_{k=1}^{m} \big(v_k - v_{k-1}\big)\, H_x\big(Y_{v_k}\big)
\;+\; \big(1 - v_m\big)\cdot \diag,
\end{equation}
where $v_1 < \cdots < v_m$ are the distinct probability values, $v_0 := 0$. The
indexing is the whole content of the identity and is easy to get wrong: since
$Y_t = \{i : p_i \ge t\}$, the level set is \emph{constant} for
$t \in (v_{k-1}, v_k]$ and equal to $Y_{v_k}$ there, so each level carries the gap
\emph{below} it, not above. The final term is the atom $Q_p$ places on the empty mask
(Remark~\ref{rem:atom}).

The \emph{exact} quadrature nodes are therefore the image's own order statistics,
weighted by their gaps, plus an atom at the empty mask; and the constants
$\alpha, 1-\alpha$ are a two-point approximation to this. Adapting \emph{nodes} per
image is legitimate --- we still estimate the same
functional for every image, so the cross-image ranking survives and only quadrature
error (a) falls. Normalizing the \emph{score} per image would not be: it would divide
out the very signal we rank by. And the weights must remain the gaps: equal weights on
quantile nodes silently swaps the posterior for one whose marginals are the image's
own empirical CDF $F_x(p_i)$, breaking Remark~\ref{rem:classical}(i).

\paragraph{A negative result: rank-anchoring the conservative threshold fails.}
Equation~\eqref{eq:orderstat} suggests an obvious repair --- replace the constant
conservative threshold with an order statistic of the class's own probabilities, so
that $\Yhi\ne\varnothing$ \emph{by construction}:
\begin{equation}\label{eq:rankanchor}
t_{\mathrm{hi}} \;=\; \min\!\Big(1-\alpha,\; \mathrm{P}_{95}\big(\{p_i : i \in \Yhat_c\}\big)\Big).
\end{equation}
It has the apparently decisive property that where the model is confident it returns
$1-\alpha$ and the score is \emph{bit-for-bit} the constant-threshold one, so it
looks like a strict information gain on the degenerate images and a no-op elsewhere.
We implemented it, and it does remove the saturation (tie fraction $0.90\to0.04$ on
a controlled synthetic population). \textbf{It also makes the ranking worse, in every
regime.} Kendall $\tau$ against the ground-truth risk falls by $0.10$--$0.15$ and
AURC rises --- under the boundary estimand \emph{and} under a convention-free
$\int_0^1 (1-\Dice)\,dt$, so this is not an artifact of the $\diag$ convention.

The reason is worth stating, because it is a trap the framework sets for itself. A
percentile of the \emph{in-mask probability values} is, by definition, an \emph{area}
quantile of the mask: $\{p \ge \mathrm{P}_{95}(p|_{\Yhat_c})\} \cap \Yhat_c$ is the
top-$5\%$-of-the-mask-by-area core (measured $|\Yhi|/|\Yhat_c|$: median $0.055$). So
\eqref{eq:rankanchor} is not an anchor in probability at all --- it is the
fixed-\emph{area} rule, and the band then reads the distance from the skirt to a
core of pinned relative area, whose inner term is pure object geometry. That is the
shape-statistic collapse the whole design exists to avoid, arriving by the back door.
It additionally breaks the nesting \eqref{eq:nesting} on normalized multi-class maps
(a softmax argmax winner need only clear $1/C$, so $t_{\mathrm{hi}}$ can fall below
$\tfrac12$ and $\Yhi\not\subseteq\Yhat_c$), where it can invert the score outright.

We report \eqref{eq:rankanchor} as an ablation, not a method.

\paragraph{A second negative result: importance sampling by the probability histogram
fails too, and for a reason that unifies both failures.} \eqref{eq:orderstat} also
suggests placing the nodes at the quantiles of the image's own probability
\emph{histogram} and weighting them by the gaps they span --- textbook importance
sampling, and, unlike \eqref{eq:rankanchor}, provably estimand-preserving, since the
gap weights are exactly the ones that leave $\int_0^1 H_x$ invariant. We implemented
it. \textbf{It is worse than the equispaced midpoint rule at every $M<16$}: on a
smoothed random field, $|r_M - \int_0^1 H_x|$ is $19.5$ vs $1.1$ at $M=2$ and $5.6$ vs
$1.4$ at $M=8$.

The mechanism is the point, and it is Observation~\ref{obs:vshape}. $H_x$ is
empirically V-shaped, with a \emph{minimum} near $m_c$ --- and, for a binary map with
\emph{no pixel tied at exactly} $\tfrac12$, $Y_{1/2} = \Yhat$ and $H_x(\tfrac12)=0$
exactly (Prop.~\ref{prop:floor}(i)). It then rises toward \emph{both} tails on the maps
we measured (one binary map: $H_x = 0.0$ at $t=\tfrac12$, $19.0$ at $t=0.25$, $78.7$ at
$t=0.98$), though neither unimodality nor $H_x(1)=\diag$ is guaranteed in general
(Remark~\ref{rem:vshapefails}, Prop.~\ref{prop:floor}(iii)). \emph{The integral's mass
lives in the tails; the neighbourhood of the minimum contributes almost nothing.} But the
probability \emph{values} concentrate near the vertex --- that is what makes it the
vertex. Placing nodes at their quantiles therefore puts them precisely where the
integrand is smallest: for the map above, the $M=2$ quantile nodes land at $0.359$ and
$0.597$, hugging the minimum, while the midpoint rule's $0.25$ and $0.75$ sit where the
mass is.

This is the same error as \eqref{eq:rankanchor}, which also pulls the conservative
node \emph{toward} the probability mass. Both failures share one lesson:

\begin{quote}
\emph{The correct importance density is $\propto |H_x(t)|$, which is tail-heavy. One
should oversample \textbf{away} from where the probability mass sits, not toward it.
Every distribution-aware node rule we tried does the opposite.}
\end{quote}

It also retro-justifies the constant threshold's nodes $(\alpha, 1-\alpha) = (0.1,
0.9)$: they are in the tails, where $H_x$ is large. That was not by design, but it was
not luck either --- it is the only region that carries signal.

\subsection{Deriving nodes from the integrand --- and why the derivation does not
settle them}\label{sec:derivednodes}

One can try to go further than ``take the nodes from a quadrature rule'' and derive
them from the shape of $H_x$. We set out that derivation, and then show --- twice ---
that it does not license the nodes it produces: minimizing the Koksma bound does not
minimize the error (Remark~\ref{rem:minimaxfallacy}), and the ``tuning buys nothing''
result holds on one synthetic generator and fails on a second. \emph{Node placement is
settled empirically (\S\ref{sec:openitems}), not here.} An earlier version of this
subsection was titled ``the nodes \emph{are} derivable, and tuning cannot beat them'';
both halves of that are retracted below.

Two facts about $H_x$ \emph{are} provable, and one regularity is only empirical. The
distinction matters, and we have already got it wrong once.

\begin{proposition}[The tightest enclosing level set, and a floor for it]\label{prop:floor}
Write $m_c := \min\{\, p_{ci} : i \in \Yhat_c \,\}$, the smallest probability the map
attains inside its own prediction. Then
\begin{enumerate}
\item $Y_t \supseteq \Yhat_c \iff t \le m_c$, so $Y_{m_c}$ is the \emph{tightest} level
set containing the prediction. Note the two sets are defined by inequalities of
\emph{different strictness}: the argmax breaks ties toward the lowest class index, so
$\Yhat_c$ excludes a pixel where $p_c$ merely ties the leader, while
$Y_t = \{i : p_{ci} \ge t\}$ includes it. In the binary case this means
$\Yhat = \{p > \tfrac12\}$ \emph{strictly}, and $Y_{1/2} = \Yhat$ --- hence
$H_x(\tfrac12)=0$ --- \textbf{only if no pixel attains exactly $\tfrac12$}. On a map
of $121$ tied pixels we measure $|\Yhat| = 9$ against $|Y_{1/2}| = 121$ and
$H_x(\tfrac12) = 5.66$, not $0$;
\item $m_c \ge \phi$, the \textbf{argmax floor}, where
$\phi = 1/C$ for a normalized (softmax) map --- a winner is the largest of $C$ numbers
summing to $1$ --- and $\phi = \tfrac12$ for a binary map and for CLIPSeg's
$\mathrm{bg} = 1 - \max_c p_c$, where a foreground class wins only if $p_c > 1-p_c$;
\item $H_x(1) = L(\varnothing,\Yhat_c) = \diag$ from \eqref{eq:bounded} --- \emph{but
only when $\max_i p_{ci} < 1$.} A confident softmax saturates to exactly $1$ in
floating point, and then $Y_1 = \{p_c \ge 1\}$ is \emph{non-empty} and $H_x(1)$ is not
the diagonal at all: on a map with a saturated core we measure $H_x(1) = 0.000$ against
$\diag = 28.28$. The atom of Remark~\ref{rem:atom} degrades gracefully --- its mass
$1-\max_i p_{ci}$ simply vanishes --- but the endpoint value does not, and any argument
that assumes $H_x(1)=\diag$ (we had one, and have withdrawn it) is unsound on
saturating maps.
\end{enumerate}
\end{proposition}

\begin{observation}[$H_x$ is empirically V-shaped, with its vertex near $m_c$]\label{obs:vshape}
On realistic maps $H_x$ falls then rises, with a minimum near $t = m_c$. Measured
monotonicity violations: $0\%$ across binary, softmax ($C=5,21$) and CLIPSeg-style
maps. \textbf{This is an empirical regularity, not a theorem}, and
Remark~\ref{rem:vshapefails} gives counterexamples.
\end{observation}

\begin{center}
\small
\begin{tabular}{lcccc}
\toprule
construction & floor $\phi$ & $m_c$ (min in-mask) & measured $t^\star$ & better predictor \\
\midrule
binary                & $0.500$ & $0.500$ & $0.500$ & (tie, exact) \\
softmax $C=5$         & $0.200$ & $0.206$ & $0.223$ & $m_c$ \\
softmax $C=21$ (VOC)  & $0.048$ & $0.055$ & $0.060$ & $m_c$ \\
CLIPSeg-style $C=21$  & $0.500$ & $0.533$ & $0.550$ & $m_c$ ($3\times$ closer) \\
\bottomrule
\end{tabular}
\end{center}

\begin{remark}[Where we were wrong, twice, and what an independent check caught]\label{rem:vshapefails}
An earlier draft asserted $H_x(\tfrac12)=0$ unconditionally (true only for $C=2$: under
a $C$-way argmax a class can win with $p_c$ well below $\tfrac12$, so
$\{p_c\ge\tfrac12\}$ is a strict \emph{subset} of $\Yhat_c$ --- $140$ pixels against
$298$ on a $C{=}5$ map, giving $H_x(\tfrac12)=2.83$). A second draft then asserted that
the vertex \emph{is} the argmax floor $\phi$. An independent adversarial check refuted
that too, and the refutation is sharp:
\begin{enumerate}
\item \textbf{$\phi$ is a lower bound on $m_c$, not the vertex.} On a $3$-way softmax
whose winner never descends below $0.40$, $H_x(1/3) = 5.66$ while $H_x(0.35)=0$: the
floor is not attained, so it is not where the minimum sits. \emph{We had the evidence
and rationalized it away} --- our measured $t^\star$ sat consistently \emph{above}
$\phi$ ($0.060$ vs $0.048$; $0.222$ vs $0.200$), which we explained as ``trimming
higher matches better'' rather than as the floor simply not being the vertex.
\item \textbf{$H_x(t^\star)=0 \iff C=2$ is false in both directions.} A $C{=}3$ map with
$Y_{1/3} = \Yhat_c$ gives $H_x(1/3)=0$.
\item \textbf{Unimodality is not guaranteed.} $H_x$ can fall, rise, and fall again: on a
map whose prediction is a single pixel, $H_x = 27.3,\,4.0,\,19.8,\,0.0$ at
$t = 0.01,\,0.10,\,0.30,\,0.45$.
\item \textbf{$H_x = 0$ does not imply $Y_t = \Yhat_c$.} $\HDF$ is a $95$th percentile,
so a one-pixel disagreement is truncated away. This is the same non-metric behaviour
that breaks Prop.~\ref{prop:band} (Remark~\ref{rem:nobridge}); it recurs wherever we
treat $\HDF$ as if it measured set equality.
\end{enumerate}
The \emph{geometric} observation behind Remark~\ref{rem:straddle} survives all of this:
on a $21$-way softmax the vertex sits near $0.06$, so nodes at $(0.1,0.9)$ or
$(0.25,0.75)$ do both sit on the same side of it. What does \emph{not} survive is the
\emph{inference} we drew from it. Remark~\ref{rem:straddle} conjectured that this
non-straddling explains our largest empirical result, and that conjecture is now
refuted twice over --- see the remark itself, which we have rewritten as the refutation.
The observation is real; the explanation it was made to carry is withdrawn. And there is
in any case \emph{no closed form} for $t^\star$ --- $m_c$ is the best cheap estimate we
have, and it is an estimate.
\end{remark}

Split at the vertex and each half is (empirically) monotone, hence of bounded
variation. For a BV integrand the composite midpoint rule is \emph{minimax}-optimal: it
attains the minimal star-discrepancy $D^\star = 1/(2M)$ in one dimension, hence the
smallest \emph{worst-case} Koksma bound \eqref{eq:koksma}. At $M=2$ that suggests one
node at the midpoint of each half:
\begin{equation}\label{eq:derivednodes}
\boxed{\; t_1 = \tfrac{1}{2}\hat{t}, \qquad t_2 = \tfrac{1}{2}\big(\hat{t} + 1\big),
\qquad \hat{t} := m_c \;}
\end{equation}
with weights equal to the subinterval widths, so only the node \emph{placement}
adapts and the estimand $\int_0^1 H_x$ is untouched. This reduces to $(0.25, 0.75)$
exactly when $m_c = \tfrac12$ (binary). It carries \emph{no tuned parameter}, but
$\hat{t} = m_c$ is an \emph{estimate} of the vertex, not the vertex
(Remark~\ref{rem:vshapefails}), so we call the rule derived-but-approximate and test
it rather than assert it. Note Gauss--Legendre is \emph{not} the right rule here despite
being the textbook choice: its optimality assumes polynomial smoothness, and $H_x$ has
a kink at its vertex. On a V-shaped $H_x \propto |2t-1|$ the midpoint rule is exact
while Gauss--Legendre overestimates by $15.5\%$; measured on real binary maps, midpoint
errs by $18.4\%$ against Gauss--Legendre's $20.8\%$.

\begin{remark}[A hypothesis of ours, refuted twice --- the explanandum stands, the
explanation does not]\label{rem:straddle}
\emph{We state this as we believed it, then refute it, because the refutation is the
content. The explanandum is not in doubt: the $5.6$--$17.6\%$ reproduces exactly on
clean data (\S\ref{sec:quadrature}). What dies is our explanation of it.}

\textbf{The hypothesis.} On a $21$-way softmax the vertex sits at $t^\star\approx0.05$.
The deployed band's nodes $(0.1, 0.9)$ and the binary-derived $(0.25, 0.75)$ therefore
put \emph{both} nodes on the \emph{same side} of it: the quadrature never straddles the
minimum, and the entire left branch of the V goes unsampled. That is exactly the regime
in which dense $M{=}32$ beats $M{=}2$ by $5.6$--$17.6\%$ --- on VOC, and only on VOC
(\S\ref{sec:quadrature}) --- while the ladder is flat on binary Pet, where $(0.25,0.75)$
straddles correctly. The reading would be that the dense rule wins on multi-class not
because $32$ nodes are needed, but because $2$ badly-placed ones sample only half the
integrand.

\textbf{First refutation: the gain is largest exactly where the mechanism is switched
off.} The mechanism can only operate where the band fails to straddle the vertex, and
that is a property of the \emph{map}, not of the dataset. By Prop.~\ref{prop:floor}(ii)
only the two genuine $21$-way softmax conditions have a low floor
($\phi = 1/21 = 0.048$); the two CLIPSeg-VOC conditions set
$\mathrm{bg} = 1-\max_c p_c$, which forces $\phi = \tfrac12$, so their vertex sits above
$\tfrac12$ and the band's $(0.1,0.9)$ \emph{already straddles it}. (Confirmed per
condition: the $5$th percentile of $\max_c p_c$ dips below $\tfrac12$ on exactly those
two softmax conditions --- $0.359$ and $0.309$ --- and never on the other six, where it
is $\ge 0.501$.) So the hypothesis makes a sharp prediction about \emph{which} VOC cells
should gain, and it gets it backwards:

\begin{center}
\small
\begin{tabular}{lccc}
\toprule
condition & $\phi$ & band straddles? & band $\to M{=}32$ gain (\% of band AURC, $95\%$ CI) \\
\midrule
\texttt{clipseg-general}, VOC    & $\tfrac12$  & \textbf{yes} & $\mathbf{+17.6}$ \; $[15.0,\ 19.9]$ \\
\texttt{clipseg-target}, VOC     & $\tfrac12$  & \textbf{yes} & $\mathbf{+14.7}$ \; $[12.3,\ 17.2]$ \\
\texttt{deeplabv3-external}, VOC & $0.048$ & no  & $+6.9$ \; $[1.9,\ 11.4]$ \\
\texttt{deeplabv3-target}, VOC   & $0.048$ & no  & $+5.6$ \; $[1.3,\ 10.0]$ \\
\bottomrule
\end{tabular}
\end{center}

The two cells where straddling \emph{provably cannot operate} gain two to three times as
much as the two where it could, and the confidence intervals \textbf{do not overlap}.
This is not a null result; it is the wrong sign. Note also what this does to the ``only
on VOC'' framing that made the hypothesis attractive: ``VOC vs Pet'' and ``low floor vs
$\phi=\tfrac12$'' are different partitions of the eight conditions, and we had been
reading a result that follows the first as though it followed the second.

\textbf{Second refutation: the discriminating experiment fails.} \texttt{vtx2} ---
\eqref{eq:derivednodes} at $\hat t = m_c$, which straddles the vertex \emph{by
construction} --- was built to test precisely this, and it \textbf{loses}
(\S\ref{sec:openitems}, item 4): $4/8$ against the deployed band, worst of the three
rules compared, and on the only two conditions where the mechanism can operate it is
\emph{worse than the band it was meant to repair}.

\emph{Caveat, stated because it bounds the refutation.} \texttt{vtx2} is not a pure
straddle toggle: it also moves the upper node to $\approx 0.53$ --- off the right tail,
where \S\ref{sec:threshold} shows the integral's mass lives --- and reweights to
$\approx(0.055, 0.945)$. So what this second argument refutes is that
\eqref{eq:derivednodes}'s straddle explains the ladder, \emph{not} that no straddling
rule whatever could help. The first argument needs no such caveat, which is why we lead
with it.

\textbf{What replaces it.} \S\ref{sec:quadrature} decomposes the $5.6$--$17.6\%$ into
legs that move one thing each, and the answer is not node placement on three of the four
cells --- nor, on the largest of them, node count.
\end{remark}

\begin{remark}[Minimax-optimal is not optimal, and we leaned on the difference]\label{rem:minimaxfallacy}
The step from ``the midpoint rule minimizes the worst-case bound'' to ``the midpoint
rule is the right node set'' is a \textbf{non-sequitur}, and an independent adversarial
check produced the counterexample. On a step integrand with $\int_0^1 H = 1.5$:

\begin{center}
\small
\begin{tabular}{lccc}
\toprule
node set & star-discrepancy $D^\star$ & Koksma bound $V\!\cdot\!D^\star$ & \textbf{actual error} \\
\midrule
midpoint $(0.25, 0.75)$ & $\mathbf{0.25}$ (minimal) & $\mathbf{1.5}$ (minimal) & $0.5$ \\
a re-tuned pair         & $0.35$ & $2.1$ & $\mathbf{0.0}$ (exact) \\
\bottomrule
\end{tabular}
\end{center}

A node set with a \emph{strictly worse} bound achieves a \emph{strictly better} error.
Minimizing a worst-case bound says nothing about the error on the integrand one
actually has. \textbf{This is the second time the same fallacy appears in this
section}: the $\sqrt{V_L/V_R}$ allocation below also provably minimizes the Koksma
bound and provably fails to reduce the error.

The same check also computes $D^\star$ for the vertex-aware nodes
\eqref{eq:derivednodes} --- $0.400$ at $\hat t = 0.20$ and $0.475$ at $\hat t = 0.05$,
both \emph{worse} than the midpoint's $0.25$ --- so by the minimax criterion those nodes
are inferior. \emph{We previously offered this as a second instance of the two criteria
disagreeing}, on the grounds that the vertex-aware nodes nevertheless measured
\emph{better} on multi-class maps. \textbf{That half is retracted.} On the real
multi-class conditions they measure \emph{worse}: \texttt{vtx2} loses $4/8$ to the band
it was built to replace and $6/8$ head-to-head against plain $(0.25,0.75)$
(\S\ref{sec:openitems}, item 4). On this integrand the two criteria happen to agree, and
the minimax criterion happens to be right.

That agreement is a coincidence and rescues nothing above. The step integrand remains a
\emph{proof} that a strictly worse bound can achieve a strictly better error, and one
family of integrands on which the bound's ordering turns out correct does not repair an
inference that is invalid in general. Minimizing a worst-case bound tells you about the
worst case; it does not tell you which node set to use on the integrand you have.

\emph{We therefore withdraw any claim that node placement is settled by theory.} What
theory gives is the $O(1/M)$ rate \eqref{eq:koksma}, and that is all. Where to put a
finite number of nodes is an empirical question, and \S\ref{sec:openitems} treats it as
one.
\end{remark}

\paragraph{On one synthetic family, tuning bought nothing --- and on another it bought
a great deal.} We grid-searched the \emph{oracle} $\alpha$ --- tuned on the very maps
it is scored on, hence an upper bound on what any tuning procedure could achieve ---
against the exact integral:

\begin{center}
\small
\begin{tabular}{lccc}
\toprule
rule (\textbf{binary maps}, $t^\star = \tfrac12$) & free parameters & mean error & nodes \\
\midrule
\textbf{midpoint $M{=}2$ (derived)} & \textbf{0} & \textbf{16.1\%} & $0.25, 0.75$ \\
symmetric 2-point, \emph{oracle} $\alpha$ & 1 & $16.1\%$ & $0.25, 0.75$ \\
asymmetric 2-point, \emph{oracle} & 2 & $16.0\%$ & $0.20, 0.75$ \\
asymmetric + free weights, \emph{oracle} & 3 & $12.1\%$ & $0.20, 0.85$ \\
\midrule
\textbf{midpoint $M{=}4$ (derived)} & \textbf{0} & \textbf{6.3\%} & --- \\
\bottomrule
\end{tabular}
\end{center}

\textbf{This table does not generalize, and we report it with its refutation
attached.} It is measured on \emph{one} synthetic family of binary maps. A second
family --- differing only in how the probability maps were generated (value spread
$\mathrm{sd}=0.065$ against $0.200$) --- puts the oracle at $(0.50, 0.78)$ with $14.4\%$
while $(0.25,0.75)$ gives $41.0\%$. \emph{Two synthetic generators, two different
answers about the optimal nodes.} The optimum depends on the distribution of
probability \emph{values}, which no theory here pins down and for which we have no
validated generator. We therefore do not claim a settled node placement, and the real
probability maps must decide it (\S\ref{sec:openitems}).

Within the one family above:
\begin{enumerate}
\item \textbf{The oracle's $\alpha$ coincided with the midpoint node.} Grid search
recovered $0.25/0.75$ and matched its error to the digit. \emph{This is a coincidence
of that generator, not a theorem} --- Remark~\ref{rem:minimaxfallacy} shows the minimax
argument cannot license it, and the second generator contradicts it. We had written it
up as ``tuning buys nothing''; that was wrong.
\item \textbf{Breaking the symmetry $t_{\mathrm{lo}}+t_{\mathrm{hi}}=1$ is not worth
it.} A second free parameter buys $0.1$ points ($16.1\to16.0$). We had argued from the
asymmetry of $H_x$ (its right tail is heavier, because the empty level set saturates
at $\diag$) that an asymmetric rule should help; it does not. The $M=2$ error is
dominated by the V's curvature, not by its asymmetry. Only freeing the \emph{weights}
helps ($\to 12.1\%$), and even that is beaten outright by the zero-parameter $M=4$.

\emph{Note the symmetry that \eqref{eq:derivednodes} \emph{does} break, and for a
reason: on a multi-class softmax $m_c$ sits far below $\tfrac12$ (measured $0.055$ on a
$21$-way map), so the derived nodes are strongly asymmetric ($\approx 0.028, 0.528$ at
$C=21$). The finding above is that asymmetry chosen by \emph{tuning} is worthless;
asymmetry \emph{derived from $m_c$} --- an \emph{estimate} of the vertex, not the argmax
floor $\phi$, which is only a lower bound on $m_c$ (Prop.~\ref{prop:floor}(ii),
Remark~\ref{rem:vshapefails}) --- is a different proposition, and it was the one
\S\ref{sec:openitems} set out to test. \emph{It has since been tested, and it is worthless
too}: the derived asymmetric nodes lose to the symmetric $(0.25,0.75)$ they were meant to
improve on, $6/8$ head-to-head on the real conditions (\S\ref{sec:openitems}, item 4). So
the honest generalization of this table is not ``tuned asymmetry is worthless while
derived asymmetry is untested'' but simply \emph{asymmetry has not helped us, however we
chose it}. An earlier version of this note said ``the vertex is at
$1/C$'' and quoted nodes $(0.024, 0.524)$ computed from $\phi = 1/21$; that is the
retracted floor rule, and the figures were a fossil of it.}
\item \textbf{Add nodes, do not tune them.} The zero-parameter $M=4$ rule ($6.3\%$)
halves the error of the three-parameter oracle at $M=2$ ($12.1\%$). Nodes cost a
distance transform; hyperparameters cost a validation split, a tuning protocol, and a
reviewer's trust.
\end{enumerate}

\begin{remark}[A derivation of ours that did \emph{not} survive]
Koksma applied per half gives $\varepsilon_L \le V_L/2M_L$ and
$\varepsilon_R \le V_R/2M_R$, so minimizing the total error bound under
$M_L + M_R = M$ prescribes the allocation $M_L/M_R = \sqrt{V_L/V_R}$ --- and both
variations are free or nearly so ($V_R = \diag$ a priori; $V_L = \HDF(\text{image},
\Yhat)$, one EDT). Measured $V_L/V_R$ has median $0.243$, so the rule sends $33\%$ of
the nodes to the low-threshold half. \emph{It does not help}: against uniform
allocation at equal $M$ it wins twice, ties twice and loses once, with margins under
one point. Minimizing a worst-case bound is not the same as minimizing the error, and
Koksma's is loose here. We report it because the derivation looks compelling and is
wrong, which is the third time in this section that a principled-looking adaptive
threshold rule has failed empirically.
\end{remark}

\paragraph{Why the band cannot inherit any of this.} The above prescribes nodes for
$r_M$, the quadrature. It says nothing about the band, and cannot: \emph{the band
width is not a quadrature, so its $\alpha$ is not a node.} The band is a surrogate
(Remark~\ref{rem:nobridge}); its $\alpha$ is a trimming level with no theoretical
home, which is exactly why the empirically optimal $\alpha$ for the band scatters over
$0.01$--$0.3$ with no pattern while the optimal node for $r_2$ sits immovably at
$0.25$.

Deploy $r_M$ instead of the band, and the problem dissolves: the nodes are
\emph{prescribed by the quadrature rule}, and there is \textbf{no free
hyperparameter at all} --- Gauss--Legendre gives $t = \tfrac12 \mp \tfrac{1}{2\sqrt3}
= 0.2113, 0.7887$ at $M=2$; the midpoint rule gives $(m-\tfrac12)/M$. And it is not a
sacrifice: the zero-hyperparameter $r_2$ \emph{beats the band with an oracle-tuned}
$\alpha$ on the CLIPSeg cells (VOC, penalized boundary: $153.0$ vs $178.8$; and
$91.0$ vs $98.9$) and ties elsewhere.

\begin{remark}[Deploying $r_M$ closes the paper's one theoretical gap]
Remark~\ref{rem:nobridge} is the only place where our theory does not reach our
practice: $\HDF$ is a percentile, not a metric, so no inequality connects the band
width to the $r_2$ it stands in for, and the bridge rests on a measured
$\rho\in[0.974,0.992]$. \emph{If we deploy $r_2$, there is nothing to bridge} --- the
deployed score \emph{is} the estimator the theory derives. The chain becomes
\begin{equation*}
\underbrace{C^\star_L = -\Ex[L]}_{\text{definition}}
\;\to\;
\underbrace{\textstyle\int_0^1 L(Y_t,\Yhat)\,dt}_{\text{exact under } Q_p}
\;\to\;
\underbrace{r_M,\ \text{nodes from the rule}}_{O(1/M) \text{ by Koksma}}
\;\to\;
\underbrace{M=2}_{\text{deployable}} ,
\end{equation*}
with no step tuned and no step unproved. We recommend it --- and, since the first draft
of this remark, the recommendation has stopped costing anything: at the band's \emph{own}
nodes $(0.1,0.9)$, $r_2$ ranks strictly better than the band on all eight conditions
under all three risks, $24/24$ (\S\ref{sec:openitems}, item 2). Closing the theoretical
gap and improving the score turn out to be the same edit.
\end{remark}

\paragraph{Selecting $\alpha$ label-free.} Three rules, all computable on the
\emph{unlabelled} test set, preserving the no-tuning-data property:
\begin{enumerate}
\item \textbf{Degeneracy budget.} Choose the smallest $\alpha$ with
$\hat\Prb(\Yhi=\varnothing)\le\delta$ (say $5\%$). Directly targets the mechanism
above; on clipseg-general it forces $\alpha \gg 0.1$.
\item \textbf{Pooled quantile rule.} Fix $1-\alpha$ at a quantile of the model's
\emph{pooled} foreground-probability distribution over the unlabelled split, rather
than at a fixed probability. This adapts to diffuse vs.\ sharp models while keeping
$\alpha$ \emph{constant across images} --- which is what distinguishes it from the
failed per-image rule \eqref{eq:rankanchor}: a global threshold cannot encode any
individual image's shape.
\item \textbf{Conformal calibration.} Calibrate $\alpha$ so that
$\Prb(\Yhi \subseteq Y^* \subseteq \Ylo) \ge 1-\delta$. Nested level sets are the
native object of Conformal Risk Control \cite{crc}, which gives $\alpha$ a
distribution-free meaning and turns the band into a genuine confidence region for
the boundary. This needs labels once, and is the most principled option.
\end{enumerate}

\begin{remark}[Temperature is the same knob]
A monotone recalibration of $p$ leaves the argmax and every $p$-ranking unchanged,
so it is a no-op for aMSP-style baselines --- but it \emph{moves the level sets} and
therefore changes our band. For this score, temperature scaling and $\alpha$
selection are the same degree of freedom. This is both a confound to control and a
point of difference from the aggregation baselines.
\end{remark}

\section{Experiments}\label{sec:experiments}

\subsection{Protocol}

\paragraph{Conditions and data.} A $2\times2$ design over model family and
adaptation, stressing scores across sharply and diffusely calibrated maps: CLIPSeg
(ViT-B/16) zero-shot and decoder-fine-tuned; DeepLabV3-ResNet50 with the COCO
checkpoint and fine-tuned end to end. Datasets: Oxford-IIIT Pet (binary,
$3{,}669$ test images) and PASCAL VOC 2012 ($21$ classes, $1{,}449$ val images).
Probability maps are bilinearly upsampled to native resolution, so distances are in
original-image pixels. This gives the \textbf{eight} condition $\times$ dataset
cells of Tables~\ref{tab:overlap}--\ref{tab:boundary}; all counts below are out of
$8$.

\paragraph{Baselines ($16$).} Mean / 5th-percentile / worst-patch maximum softmax
probability, mean margin, Shannon entropy (global and foreground-restricted),
low-confidence fraction, tail-mean uncertainty \cite{values}, MMMC
\cite{holder}, DOCTOR's collision probability \cite{doctor}, Generalized Entropy
\cite{gen}, boundary and interior entropy \cite{metaseg,zenk}, Boundary Uncertainty
Concentration \cite{buc}, Soft Dice Confidence \cite{sdc}, and the SAM stability
score \cite{sam} (Remark~\ref{rem:sam}). We additionally evaluate a \emph{combined}
score, the rank-average of the distance band and SDC.

\emph{The combined score is \textbf{transductive}, and we mark it wherever it appears.}
Its two rank transforms are taken across the evaluation split, so its value --- and, for
a small fraction of pairs, its \emph{ordering} --- depends on which other images are
present, and scored one image at a time it is constant. It is therefore \emph{not} a
per-image functional $g(x)$ and falls outside the post-hoc, image-level,
\emph{single-pass} regime this paper defines in \S\ref{sec:setup}. It leaks no ground
truth, so its AURC is not inflated; but it \textbf{cannot be deployed as written}, and
no recommendation should be read from it. The deployable version freezes the two
empirical CDFs on a held-out split; we have not measured that. The code marks it
(\texttt{TRANSDUCTIVE\_SCORES} in \texttt{scripts/analyze\_selective.py}) and so, from
here on, do we.

\paragraph{Metric.} AURC, with \emph{random} and \emph{oracle} bracketing every
achievable score. Significance: a paired image bootstrap ($1000$ resamples) of the
band against SDC.

\subsection{Results}\label{sec:results}

\begin{table}[t]
\centering
\small
\caption{AURC under \textbf{overlap risk} ($1-\mathrm{mIoU}$), lower is better;
best in bold. Band-HD95 is \eqref{eq:bandsym}; Band-Dice and SAM-stab are the area
readouts \eqref{eq:area} (Remark~\ref{rem:sam}); bi-level scores at $\alpha=0.1$,
symmetric, present-class mean. Detection failures are ranked least-confident
(\S\ref{sec:multiclass}). \textbf{Combined is \emph{transductive}} --- its ranks are
taken across the split, so it is not a per-image score and is not deployable as written
(\S\ref{sec:experiments}).}
\label{tab:overlap}
\begin{tabular}{lccccccccc}
\toprule
Condition & Band-HD95 & Band-Dice & SAM-stab & SDC & Combined & aMSP & aNE & BUC & Oracle \\
\midrule
clipseg-general, Pet     & 0.0176 & 0.0191 & 0.0191 & 0.0204 & \textbf{0.0171} & 0.0178 & 0.0192 & 0.0179 & 0.0057 \\
clipseg-general, VOC     & 0.2920 & 0.2882 & 0.2910 & \textbf{0.2375} & 0.2444 & 0.3321 & 0.3314 & 0.3263 & 0.1940 \\
clipseg-target, Pet      & \textbf{0.0047} & 0.0126 & 0.0126 & 0.0126 & 0.0059 & 0.0076 & 0.0072 & 0.0049 & 0.0013 \\
clipseg-target, VOC      & 0.1777 & 0.1781 & 0.1807 & 0.1627 & \textbf{0.1601} & 0.2236 & 0.2260 & 0.2202 & 0.1285 \\
deeplabv3-external, Pet  & \textbf{0.0118} & 0.0179 & 0.0179 & 0.0176 & 0.0124 & 0.0211 & 0.0197 & 0.0131 & 0.0036 \\
deeplabv3-external, VOC  & 0.1018 & 0.1072 & 0.1091 & 0.1072 & \textbf{0.0993} & 0.1523 & 0.1451 & 0.1327 & 0.0670 \\
deeplabv3-target, Pet    & 0.0060 & 0.0145 & 0.0145 & 0.0140 & 0.0072 & 0.0082 & 0.0074 & \textbf{0.0054} & 0.0016 \\
deeplabv3-target, VOC    & 0.0912 & 0.0949 & 0.0962 & 0.0950 & \textbf{0.0885} & 0.1340 & 0.1297 & 0.1185 & 0.0617 \\
\bottomrule
\end{tabular}
\end{table}

\begin{table}[t]
\centering
\small
\caption{AURC under \textbf{boundary risk} ($\HDF^{\mathrm{drop}}$, pixels), lower
is better. BdEnt = boundary entropy. Both spatially aware baselines (BUC, BdEnt) are
included because both beat the band in one cell. \textbf{Combined is
\emph{transductive}} and not deployable as written (\S\ref{sec:experiments}).}
\label{tab:boundary}
\begin{tabular}{lccccccccc}
\toprule
Condition & Band-HD95 & Band-Dice & SAM-stab & SDC & Combined & aMSP & BUC & BdEnt & Oracle \\
\midrule
clipseg-general, Pet     & 9.068 & 10.015 & 10.015 & 9.387 & 8.749 & 9.373 & 9.567 & \textbf{8.540} & 2.771 \\
clipseg-general, VOC     & 26.068 & 27.131 & 27.143 & 26.947 & \textbf{25.093} & 31.703 & 31.547 & 39.846 & 12.728 \\
clipseg-target, Pet      & \textbf{2.012} & 3.416 & 3.416 & 3.362 & 2.212 & 2.555 & 2.117 & 2.309 & 0.487 \\
clipseg-target, VOC      & 17.597 & 17.774 & 17.610 & 18.530 & \textbf{17.543} & 22.857 & 22.697 & 22.101 & 9.022 \\
deeplabv3-external, Pet  & \textbf{5.387} & 7.257 & 7.257 & 7.235 & 5.492 & 7.929 & 6.874 & 6.887 & 1.771 \\
deeplabv3-external, VOC  & 13.479 & 14.388 & 14.311 & 14.592 & \textbf{13.471} & 19.366 & 17.909 & 16.538 & 7.293 \\
deeplabv3-target, Pet    & 2.537 & 4.152 & 4.152 & 4.081 & 2.717 & 2.996 & \textbf{2.409} & 3.175 & 0.673 \\
deeplabv3-target, VOC    & \textbf{13.244} & 15.070 & 14.902 & 15.298 & 13.624 & 18.631 & 17.280 & 15.209 & 7.098 \\
\bottomrule
\end{tabular}
\end{table}

\begin{table}[t]
\centering
\small
\caption{AURC under the \textbf{penalized boundary risk} $\HDF^{\mathrm{pen}}$
(detection failures charged the image diagonal, so every risk is evaluated on the
same image set), lower is better. Compare Table~\ref{tab:boundary}: the
distance-vs-area ablation is \emph{unchanged} ($8/8$), but the band's win over SDC
falls from $8/8$ to $6/8$, losing exactly the two CLIPSeg-VOC cells --- the same two
it loses under overlap risk. \textbf{Combined is \emph{transductive}} and not deployable
as written (\S\ref{sec:experiments}).
\emph{An earlier version of this table carried a warning that its
\texttt{clipseg-target}/VOC row was stale, disagreed with every other report of the
same cell in this paper, and must not be cited until all three tables were
regenerated from a single post-fix evaluation. The regeneration has happened. The
warning was right --- that row was populated from a different run, and five of the
eight rows moved --- and its prediction was also right: the disagreement flipped no
cell, and both the $6/8$ count and the $8/8$ distance-vs-area count stand exactly as
they were reported under embargo. Every row here is now the single clean evaluation,
and the cell agrees with the \eqref{eq:mband} table and the $M$-ladder at band
$101.6$, SDC $87.5$.}}
\label{tab:penalized}
\begin{tabular}{lccccccccc}
\toprule
Condition & Band-HD95 & Band-Dice & SAM-stab & SDC & Combined & aMSP & BUC & BdEnt & Oracle \\
\midrule
clipseg-general, Pet     & \textbf{9.62} & 11.41 & 11.41 & 11.55 & 9.63 & 11.13 & 11.02 & 17.26 & 2.82 \\
clipseg-general, VOC     & 182.91 & 185.03 & 187.30 & \textbf{135.67} & 145.59 & 209.74 & 203.44 & 197.26 & 102.15 \\
clipseg-target, Pet      & \textbf{2.31} & 7.50 & 7.50 & 7.48 & 3.19 & 4.59 & 2.71 & 2.35 & 0.50 \\
clipseg-target, VOC      & 101.65 & 104.71 & 106.97 & \textbf{87.45} & 87.46 & 135.89 & 131.51 & 117.62 & 62.82 \\
deeplabv3-external, Pet  & \textbf{6.69} & 10.13 & 10.13 & 9.85 & 7.10 & 15.50 & 8.22 & 17.32 & 1.84 \\
deeplabv3-external, VOC  & 47.88 & 52.72 & 54.39 & 50.56 & \textbf{47.07} & 90.90 & 68.52 & 89.84 & 25.27 \\
deeplabv3-target, Pet    & 3.11 & 7.77 & 7.77 & 7.30 & 3.60 & 5.10 & \textbf{3.06} & 3.63 & 0.69 \\
deeplabv3-target, VOC    & 46.38 & 51.42 & 52.72 & 50.32 & \textbf{46.22} & 81.48 & 63.39 & 79.31 & 25.45 \\
\bottomrule
\end{tabular}
\end{table}

\paragraph{Finding 1 (core): the distance readout dominates its own area readout ---
$8/8$ under \emph{both} boundary conventions, and $7/8$ under overlap risk too.} Among
the readouts of the \emph{same} band
--- which share posterior, thresholds and aggregation, and differ only in whether
they are permutation invariant (Lemma~\ref{lem:perm}) --- the distance form wins
boundary risk in \textbf{$8/8$} cells under $\HDF^{\mathrm{drop}}$
(Table~\ref{tab:boundary}) \emph{and} $8/8$ under $\HDF^{\mathrm{pen}}$
(Table~\ref{tab:penalized}). It also wins \emph{overlap} risk in $7/8$ of the cells of
Table~\ref{tab:overlap}, losing only \texttt{clipseg-general}/VOC (and winning
\texttt{clipseg-target}/VOC by only $0.2\%$, so read that cell as a coin flip). Spatial
information helps most when the risk is spatial, and does not hurt when it is not.
\emph{An earlier version of this paragraph claimed ``an area form wins overlap risk in
$5/8$'' --- contradicted by our own Table~\ref{tab:overlap}, in which the area readout
wins exactly one cell. It was a leftover of a superseded analysis, and the corrected
count is the stronger result.} This is the cleanest test of the
thesis because the comparison is controlled, and it is the one result in the paper
that is invariant to every convention we vary --- including the two that
\S\ref{sec:results} shows can flip the SDC comparison. The distance transform also
beats its own perimeter-normalized proxy $8/8$ under \emph{every} risk, by a median
of $1.17\times$ AURC under $\HDF^{\mathrm{drop}}$ (max $1.39\times$) and
$1.24\times$ under both $\HDF^{\mathrm{pen}}$ and overlap (max $2.08\times$) --- the
geometry, not merely the band's size, carries the signal.

\paragraph{Finding 2: the band beats SDC $6/8$ under \emph{both} risks, and fails on
the \emph{same two cells} in both --- so this is a condition effect, not a risk
effect.} Under overlap risk the band beats SDC in $6/8$. Under the boundary risk it
beats SDC $8/8$ with the $\HDF^{\mathrm{drop}}$ convention
(Table~\ref{tab:boundary}) but only \textbf{$6/8$} once detection failures are
charged rather than deleted (Table~\ref{tab:penalized}) --- \emph{the $8/8$ headline
does not survive the honest convention, and we report it degraded.} In every case
the two losses are the same two cells: \texttt{clipseg-general}/VOC and
\texttt{clipseg-target}/VOC.

That coincidence is the finding. A metric-alignment story cannot explain why a
score fails on the same conditions under \emph{both} an overlap and a boundary risk.
A condition-level story can, and Finding 4 gives it.

Two of the paper's own remedies each recover part of the gap, and neither is a
threshold fix: restoring the nesting \eqref{eq:mband} closes $42$--$77\%$ of it
(Finding 4), and dense quadrature (\S\ref{sec:quadrature}) takes back
\texttt{clipseg-target}/VOC outright, lifting the count to $\mathbf{7/8}$. An earlier
draft withheld that second result while the ladder that measured it was audited for the
aggregation defect of Remark~\ref{rem:mladdercontam}; the recomputation has landed, the
defect never touched the ladder, and we state it affirmatively. The single cell that
survives everything is
\texttt{clipseg-general}/VOC ---
zero-shot CLIPSeg on the multi-class benchmark, where the detection-failure rate
($67\%$) and the nesting leak ($7.0\%$) are both at their worst.

\paragraph{Finding 3 (why the band fails where it fails): it cannot see detection
failures, and CLIPSeg-on-VOC is made of them.} The penalized risk charges each
hallucinated or wholly-missed class the image diagonal. How often that fires, and
what it does to the risk, differs by an order of magnitude across conditions:

\begin{center}
\small
\begin{tabular}{lccc}
\toprule
condition & images with a one-sided class & $\HDF^{\mathrm{pen}} / \HDF^{\mathrm{drop}}$ & band vs.\ SDC \\
\midrule
Pet (all four conditions) & $0.2$--$0.4\%$ & $1.18$--$1.28\times$ & win \\
\texttt{deeplabv3}, VOC   & $31$--$32\%$ & $3.9$--$4.1\times$ & win \\
\texttt{clipseg}, VOC     & $\mathbf{54}$--$\mathbf{67\%}$ & $\mathbf{5.8}$--$\mathbf{6.3\times}$ & \textbf{lose} \\
\bottomrule
\end{tabular}
\end{center}

On CLIPSeg-VOC a majority of images contain a class that is present on only one
side, so the penalized boundary risk there is dominated by \emph{detection} error,
not by boundary geometry. The band is a contour-displacement statistic: it measures
how far the boundary might move, and is structurally unable to say whether a class
is \emph{there at all}. SDC, which references the predicted mask's probability mass,
degrades more gracefully. This is the ``it measures blur, not error'' limitation of
\S\ref{sec:limitations}, made quantitative --- and it locates the defect precisely,
in the probability map's class structure rather than in the choice of metric.

\paragraph{Finding 4 (mechanism, confirmed): the band's precondition is a
\emph{normalized} probability map, and CLIPSeg violates it.} The bi-level
construction assumes the nesting \eqref{eq:nesting}, in particular
$\Yhi = \{p_c \ge 1-\alpha\} \subseteq \Yhat_c$. This is not a technicality --- it is
what makes the band a reading of \emph{that class's} boundary uncertainty. We
measured how often it fails, with a label-free diagnostic (the share of
conservative-set pixels lying outside their own class's prediction):

\begin{center}
\small
\begin{tabular}{llcccc}
\toprule
condition & map & leak @ $\alpha{=}0.05$ & @ $0.1$ & @ $0.2$ & band vs.\ SDC \\
\midrule
\texttt{clipseg-general}, VOC   & sigmoid & $2.2\%$ & $7.0\%$ & $15.2\%$ & \textbf{lose} \\
\texttt{clipseg-target}, VOC    & sigmoid & $2.7\%$ & $4.3\%$ & $6.7\%$  & \textbf{lose} \\
\texttt{deeplabv3-external}, VOC& softmax & $\mathbf{0.0\%}$ & $\mathbf{0.0\%}$ & $\mathbf{0.0\%}$ & win \\
\texttt{deeplabv3-target}, VOC  & softmax & $\mathbf{0.0\%}$ & $\mathbf{0.0\%}$ & $\mathbf{0.0\%}$ & win \\
\bottomrule
\end{tabular}
\end{center}

The nesting is violated on \emph{exactly} the two conditions where the band loses,
and is \emph{exactly} zero on the two where it wins. The mechanism is structural.
CLIPSeg scores VOC with $20$ \emph{independent per-prompt sigmoids} and sets
background to $1-\max_c p_c$, so semantically overlapping prompts (cat/dog,
cow/horse/sheep --- the classic VOC confusions) co-fire: two prompts both clear
$1-\alpha$ while only one wins the argmax, and the loser's conservative set escapes
its own prediction. A normalized softmax \emph{cannot} do this --- two classes cannot
both exceed $1-\alpha \ge 0.8$ on the simplex --- which is why the softmax leak is
identically zero, and why the leak grows with $\alpha$ as $1-\alpha$ falls toward
$\tfrac12$.

The failure mode is worse than a broken invariant. When the leaked pixels lie in
\emph{both} level sets, $\Ylo$ and $\Yhi$ coincide there and the band \emph{vanishes}
--- the score reports maximal confidence on precisely the images where two classes
are fighting. On a synthetic co-firing map the plain band reads $0.00$ (perfectly
confident) where the mask-intersected band reads $18.38$\,px.

\emph{This is a precondition on the method, and we state it as one: the bi-level band
requires a normalized (softmax) or binary probability map. On per-prompt sigmoid maps
--- CLIPSeg, and open-vocabulary segmenters generally --- it degrades, and the
diagnostic above says by how much, without ground truth.}

\paragraph{The repair, and what it does \emph{not} repair.} The diagnosis prescribes
its own fix: intersect the conservative set with the class's prediction,
\begin{equation}\label{eq:mband}
\Yhi \;:=\; \{p_c \ge 1-\alpha\} \;\cap\; \Yhat_c ,
\end{equation}
which restores \eqref{eq:nesting} by construction. It is a \emph{provable no-op}
wherever the leak is zero --- on a normalized map the conservative set already lies
inside the prediction --- so it cannot cost anything on softmax or binary maps, and we
confirm this to machine precision ($47.876 \to 47.876$ on \texttt{deeplabv3-external}).
On the sigmoid conditions it closes $42$--$77\%$ of the band's AURC gap to SDC:

\begin{center}
\small
\begin{tabular}{lcccc}
\toprule
condition (risk) & band & band + \eqref{eq:mband} & SDC & gap closed \\
\midrule
\texttt{clipseg-general}, VOC (overlap)  & $0.292$ & $0.269$ & $0.238$ & $42\%$ \\
\texttt{clipseg-target},  VOC (overlap)  & $0.178$ & $0.167$ & $0.163$ & $69\%$ \\
\texttt{clipseg-general}, VOC (boundary) & $182.9$ & $163.3$ & $135.7$ & $42\%$ \\
\texttt{clipseg-target},  VOC (boundary) & $101.6$ & $90.7$  & $87.5$  & $\mathbf{77\%}$ \\
\texttt{deeplabv3}, VOC (both)           & \multicolumn{2}{c}{\emph{identical (no-op)}} & --- & --- \\
\bottomrule
\end{tabular}
\end{center}

Because \eqref{eq:mband} is free where it does nothing and strictly better where it
does something, it dominates, and \emph{we recommend it as the deployed form of the
band on any non-normalized map}. To keep every number in this paper comparable to a
single, fixed score, however, the tables above and below report the
\textbf{unrepaired} band \eqref{eq:bandsym} --- the code's \texttt{DEFAULT\_BAND} ---
with \eqref{eq:mband} reported alongside it in this table and shipped as a separate
score (\texttt{neg\_mband\_width\_*}) rather than folded into the headline. Nothing
turns on the choice: \textbf{it does not flip a single cell}, and the band still loses
$2/8$ to SDC either way. The residual gap is the \emph{other}
mechanism --- detection-failure blindness --- and that one is not an implementation
flaw to be repaired but an intrinsic limit of the score class. A
contour-displacement statistic cannot be made to see whether a class exists. The two
mechanisms are separable, and only one of them is ours to fix.

\emph{We initially attributed the CLIPSeg-VOC losses to the band degeneracy of
\S\ref{sec:threshold}, and dismissed that explanation on the grounds that the
saturation rates there --- $17.7\%$ and $1.4\%$ --- were far too small to matter.
The rates are right, but the dismissal was not: those are rates of \emph{total}
degeneracy, and the rate that actually pollutes the score is the partial one, which
is $63.1\%$ and $35.9\%$ on the same two cells. ``Far too small'' is not available
to us, and the observational pattern cannot rescue it either --- ranked by partial
degeneracy the four VOC conditions are $63.1$, $35.9$, $32.1$, $27.1\%$, and the two
that lose are the two highest, so degeneracy and the leak are \textbf{collinear}
across our conditions. What decides between them is the \textbf{intervention}:
\eqref{eq:mband} repairs the leak while provably leaving saturation untouched (an
empty conservative set has an empty intersection; band and \eqref{eq:mband} are
bit-identical on every saturated image), and it closes $42$--$77\%$ of the gap. A
degeneracy explanation predicts that a leak-only repair does nothing
(Remark~\ref{rem:retraction}).}

\emph{We record a correction to our own earlier analysis.} A previous version
reported that SDC \emph{beat} the band under overlap risk on all four VOC cells and
presented this as the ``risk-alignment split'' the theory predicted. That was an
artifact of the aggregation bug in \S\ref{sec:multiclass}: total detection failures
were being ranked as the band's \emph{most} confident images. With the convention
corrected, band-vs-SDC under overlap goes from $2/8$ to $6/8$ and the split
disappears. The corrected result is the \emph{stronger} one, and Lemma~\ref{lem:perm}
never predicted a split in that direction --- it predicts only that SDC should not
lose \emph{for lack of geometry} under overlap risk, not that it should win. The
earlier framing was a post-hoc rationalization of a bug, and we flag it rather than
quietly restate it.

\paragraph{Finding 4: the area readout tracks SDC, as Lemma~\ref{lem:perm} predicts
--- except where it does not.} Band-Dice and SDC agree to within $3.5\%$ relative
AURC on most cells (e.g.\ $0.1072$ vs $0.1072$), consistent with both being
permutation-invariant functionals. They diverge on the zero-shot CLIPSeg conditions
(VOC: $0.2882$ vs $0.2375$; per-image Spearman $0.44$). They are different
functionals --- Band-Dice is a two-quantile area ratio with no reference to $\Yhat$,
SDC is a dense statistic that has one --- so we report the agreement as an empirical
observation, not as a prediction of the theory.

\paragraph{Finding 5 (negative): the spatially aware baselines are competitive.}
BUC \cite{buc} beats the band under boundary risk on fine-tuned DeepLabV3/Pet
($2.409$ vs $2.537$), and boundary entropy \cite{metaseg} beats it on
clipseg-general/Pet ($8.540$ vs $9.068$). BUC also beats \emph{SDC} under boundary
risk on three of the four Pet cells, losing only on clipseg-general/Pet ($9.567$ vs
$9.387$). \emph{We previously wrote ``on all four''; that was wrong on our own
Table~\ref{tab:boundary}, and it overstated a result that is already against us.}
Our earlier claim that ``SDC is the only genuinely
competitive external baseline'' was false by our own data. This is not damaging to
the thesis --- it is evidence \emph{for} it: the baselines that compete under
boundary risk are exactly the two that are not permutation invariant. The honest
statement is that against \emph{spatially blind} baselines the band wins $8/8$ under
boundary risk, and against spatially aware ones it wins $6/8$.

\paragraph{Finding 6: fusing the two signals is the best all-rounder --- but the fusion
we measured is \emph{not deployable}.} The rank-average of band and SDC gives the lowest
AURC in $4/8$ cells under overlap, $3/8$ under $\HDF^{\mathrm{drop}}$ and $2/8$ under
$\HDF^{\mathrm{pen}}$. \emph{An earlier version reported $4/8$ under boundary risk,
which its own Table~\ref{tab:boundary} contradicts; and the count degrades under the
honest penalized convention, like every other count in \S\ref{sec:results}.} The two
are aligned to
different risk geometries and are complementary, and that much is a real finding about
the \emph{signals}. It is \textbf{not} a recommendation of \emph{this score}: as
\S\ref{sec:experiments} states, the combined score is \textbf{transductive} --- its
ranks are taken across the split, so it is not a per-image functional and is outside our
own single-pass regime. Read Finding 6 as evidence that the two signals are
complementary, and as a call to measure the deployable, frozen-CDF fusion --- which we
have not done.

\subsection{How many thresholds? The $M$-ablation overturns our own argument}\label{sec:quadrature}

\S\ref{sec:band} argued that $M=2$ suffices, on the grounds that AURC depends only on
the ranking \eqref{eq:rank} and that under the scale family $H_x(t)\approx a_x h(t)$
a two-point rule and the dense integral induce the \emph{same} order. That argument
is correct where its premise holds, and \textbf{its premise fails on multi-class
data}.

\paragraph{The quadrature error is provably $O(1/M)$.} $H_x(t)=L(Y_t,\Yhat)$ is of
bounded variation --- it is piecewise constant, jumping only where $t$ crosses one of
the image's own probability values, and is bounded by \eqref{eq:bounded}. Koksma's
inequality then bounds term (a) of \eqref{eq:decomp} for the $M$-point midpoint rule:
\begin{equation}\label{eq:koksma}
\Big| r_M(x) - \int_0^1 H_x(t)\,dt \Big| \;\le\; \frac{V(H_x)}{2M},
\end{equation}
with $V$ the total variation. So the quadrature \emph{provably} converges to
$\Ex_{Q_p}[L]$. What no amount of $M$ can touch is term (b): $Q_p$ is a modelling
choice, and a probability map does not determine a joint law
(Remark~\ref{rem:classical}).

\begin{remark}[\textbf{The audit that demanded these numbers back, and the recomputation
that vindicated them}]\label{rem:mladdercontam}
An audit found a real aggregation defect in the ladder below \emph{after} it was
measured, and the defect pointed the same way as the result. We withheld every number in
this subsection and said so in print. They have since been recomputed on clean data.
\emph{They reproduce to the digit}, and the embargo is lifted: every AURC in the ladder
below is bit-identical between the pre-fix run and the clean regeneration. We keep the
whole episode on the page: a withheld number that comes back unchanged is evidence about
the process, not an embarrassment to be quietly deleted. (Two cells of the table below
have since had their last printed digit corrected --- $155.9$ and $45.9$, previously
typeset as $156.0$ and $46.0$ --- because the original table double-rounded them. That
is a typesetting fix, not a movement in the numbers, and the same fix was applied to
four cells of Table~\ref{tab:boundary}.)

\textbf{The defect was real.} A class can \emph{win the argmax} while its aggressive set
$\{p_c \ge \alpha\}$ is empty --- a $C$-way softmax winner needs only
$p_c \ge \phi = 1/C = 0.048$ on VOC, below $\alpha = 0.1$ --- i.e.\ it can hallucinate a
whole object at $p \approx 0.05$. The aggregation \emph{silently dropped} such a class
instead of saturating it at the diagonal, so on one image the deployed band scored
$-5.10$ where the bounded convention \eqref{eq:bounded} demands $-47.76$: a hallucinated
class was worth $42$\,px of free confidence. It is fixed and pinned by regression tests.

\textbf{The feared confound would have been lethal, and was correctly identified.} The
rules of the ladder do not all fire the same way. The argmax floor $0.048$ sits \emph{below} the
lowest nodes of \texttt{mid16} ($0.031$) and \texttt{mid32} ($0.016$) but \emph{above}
those of \texttt{gl2} ($0.211$), \texttt{mid4} ($0.125$) and \texttt{mid8} ($0.062$). So
the dense rules could see those classes and the cheap ones could not, and the ladder
would have been comparing rules computed over \emph{different sets of classes on the same
image}: the dense rule's advantage would be an artifact --- not integrating better,
merely the only rule looking. That confound is exactly confounded with
Remark~\ref{rem:straddle}'s explanation, and it was the right reason to refuse to cite
the table.

\textbf{It never materialized.} The fix landed after the affected evaluation, which
leaves a natural experiment: the pre-fix score dump and the clean regeneration can be
diffed score by score. Over all $221$ scores common to the two, the ones that moved at
all are \emph{exactly} $26$ on \texttt{deeplabv3-external}/VOC and $26$ on
\texttt{deeplabv3-target}/VOC --- \textbf{zero on the other six conditions} --- and
\textbf{every one of them is at $\alpha = 0.3$}. Three facts account for that, and each
is checkable:
\begin{enumerate}
\item \textbf{It could only ever fire on two of the eight conditions.} The defect needs
$\phi$ below the lowest threshold, and only \texttt{deeplabv3-external}/VOC and
\texttt{deeplabv3-target}/VOC are a genuine $21$-way softmax ($\phi = 0.048$). The four
Pet conditions are binary, and
both CLIPSeg-VOC conditions set $\mathrm{bg} = 1-\max_c p_c$; all six therefore have
$\phi = \tfrac12$ (Prop.~\ref{prop:floor}(ii)), which no threshold in the grid reaches.
\item \textbf{It never bit at the deployed $\alpha = 0.1$.} On those two conditions, no
present class had a maximum probability in $[0.048, 0.1)$ --- the window in which a class
wins the argmax yet has an empty aggressive set. The defect had to reach $\alpha = 0.3$
before any class fell into it, and every headline in this paper is at $\alpha = 0.1$.
\item \textbf{The ladder itself was already clean.} Every $r_M$ score's AURC delta across
the fix is exactly $0.0$ --- every condition, every $M$. The rules were never compared
over different class sets, because there was never a dropped class to differ over.
\end{enumerate}

\textbf{Verdict: the table below stands as measured.} The audit was right to demand the
recomputation --- none of the above was knowable in advance, and the confound it named
was real and would have been fatal --- and the recomputation was right to hand the
numbers back unchanged. \emph{The ``do not cite'' was correct to issue and is now
lifted.} (The distance-vs-area result never depended on any of this: both readouts share
the aggregation.) What the episode did leave behind is a scoping fact we did not have
before, and which does real work in Remark~\ref{rem:straddle}: only $2$ of our $8$
conditions have $\phi < \tfrac12$.
\end{remark}

\paragraph{Empirically, the dense rule beats the deployed band by $5.6$--$17.6\%$ on
multi-class --- but it is not the node count that buys it.} AURC
under the penalized boundary risk:

\begin{center}
\small
\begin{tabular}{lcccccc}
\toprule
condition & band ($M{=}2$) & $M{=}2$ & $M{=}4$ & $M{=}8$ & $M{=}16$ & $M{=}32$ \\
\midrule
\texttt{clipseg-general}, VOC   & $182.9$ & $153.0$ & $155.6$ & $155.9$ & $151.8$ & $\mathbf{150.8}$ \\
\texttt{clipseg-target}, VOC    & $101.6$ & $91.0$  & $89.0$  & $88.3$  & $87.3$  & $\mathbf{86.7}$ \\
\texttt{deeplabv3-external}, VOC& $47.9$  & $48.0$  & $47.3$  & $45.9$  & $45.0$  & $\mathbf{44.6}$ \\
\texttt{deeplabv3-target}, VOC  & $46.4$  & $47.6$  & $45.7$  & $45.9$  & $45.2$  & $\mathbf{43.8}$ \\
\midrule
Pet (all four)                  & \multicolumn{6}{c}{\emph{flat: spread $2.2$--$9.6$ within condition; $M{=}2$ wins $2/4$, $M{=}16$ wins $2/4$}} \\
\bottomrule
\end{tabular}
\end{center}

The dense rule wins $4/4$ multi-class cells, by $5.6$--$17.6\%$, and is flat on the
binary ones. The scale-family premise evidently holds for a single foreground class
and fails once several classes contribute to the present-class mean --- which is
exactly when $H_x$ can change shape, not just scale, from image to image.

\paragraph{What the $5.6$--$17.6\%$ actually decomposes into --- and it is not node
count.} The figure reproduces exactly on clean data ($17.56$, $14.67$, $6.90$, $5.63\%$,
all four bootstrap CIs excluding zero). What it does not do is measure what we said it
measured. Its baseline is the \emph{deployed band}, and the band is \textbf{not} the
$M{=}2$ rung of this ladder: moving from the band to $M{=}32$ changes \emph{three} things
at once.
\begin{enumerate}
\item \textbf{Readout.} The band width $\HDF(\Ylo,\Yhi)$ \eqref{eq:bandsym} becomes the
quadrature $\tfrac12[\HDF(\Ylo,\Yhat)+\HDF(\Yhi,\Yhat)]$ \eqref{eq:r2}. A different
functional, not a different number of nodes (Remark~\ref{rem:nobridge}).
\item \textbf{Placement.} The nodes $(0.1, 0.9)$ become the rule's own.
\item \textbf{Count.} $2$ becomes $32$.
\end{enumerate}
Only the third is ``more nodes''. Separating them needs an $M{=}2$ rung sharing the
band's nodes \emph{exactly}, and there is one: $r_2$ at $(\alpha, 1-\alpha) = (0.1,
0.9)$, which differs from the band only in the readout (\S\ref{sec:openitems}, item 2).
Inserting it splits the total into three legs that each move one thing ($10$k paired
image bootstrap, as \% of band AURC; $*$ marks a CI excluding zero; the legs sum to the
total):

\begin{center}
\small
\begin{tabular}{lccc}
\toprule
 & READOUT & PLACEMENT & COUNT \\
condition & band $\to r_2@0.1$ & $r_2@0.1 \to$ \texttt{gl2} & \texttt{gl2} $\to M{=}32$ \\
\midrule
\texttt{clipseg-general}, VOC    & $+7.07^*$ & $\mathbf{+9.26^*}$ & $+1.23$ \\
\texttt{clipseg-target}, VOC     & $+9.23^*$ & $+1.24$ & $+4.20^*$ \\
\texttt{deeplabv3-external}, VOC & $+0.27$ & $-0.48$ & $\mathbf{+7.10^*}$ \\
\texttt{deeplabv3-target}, VOC   & $+1.09$ & $-3.71$ & $\mathbf{+8.24^*}$ \\
\midrule
\texttt{clipseg-general}, Pet    & $+3.23^*$ & $+9.57^*$ & $\mathbf{-10.45^*}$ \\
\texttt{clipseg-target}, Pet     & $+2.32^*$ & $+2.33$ & $\mathbf{-4.35}$ \\
\texttt{deeplabv3-external}, Pet & $+1.93$ & $-0.79$ & $+5.34$ \\
\texttt{deeplabv3-target}, Pet   & $+2.23$ & $-3.91$ & $+3.08$ \\
\bottomrule
\end{tabular}
\end{center}

\textbf{The COUNT leg is significantly positive $3\times$, significantly \emph{negative}
$1\times$, and null $4\times$} (the second negative cell, \texttt{clipseg-target}/Pet at
$-4.35\%$, is borderline: its bootstrap CI upper bound sits just above zero across seeds).
On \texttt{clipseg-general}/VOC --- the \emph{top} of the
$5.6$--$17.6\%$ range, the cell the headline is carried by --- the node-count component
is $+1.23$ $[-0.77, +3.40]$, indistinguishable from zero. \emph{So ``more nodes help'' is
not what that range shows}, and the paragraph heading this subsection carried for several
drafts was wrong.

The honest reading is that the range is a \textbf{mixture} of two unrelated stories that
happen to land in the same interval. On DeepLabV3-VOC the gain is essentially all COUNT,
with READOUT and PLACEMENT null or negative. On CLIPSeg-VOC it is essentially all READOUT
$+$ PLACEMENT, with COUNT contributing nothing on the larger of the two. And on Pet,
which we reported as ``flat'', the ladder is not flat but \emph{cancelling}: a positive
READOUT and PLACEMENT against a negative COUNT (significantly so on
\texttt{clipseg-general}/Pet). Three cells' worth of
``denser quadrature integrates better'' and one cell's worth of ``the deployed nodes are
misplaced'' were being read off a single averaged range. The one leg that is positive on
all eight conditions --- and significant on four --- is READOUT, which costs nothing
(\S\ref{sec:openitems}, item 2).

\emph{The methodological trap, recorded because we walked into it.} The obvious
decomposition is the two-leg \texttt{band} $\to$ \texttt{gl2} $\to$ $M{=}32$, reading the
first leg as ``readout'' since \texttt{gl2} is the ladder's own $M{=}2$ rung. \textbf{It
is not identified.} \texttt{gl2}'s nodes are $(0.211, 0.789)$ and the band's are $(0.1,
0.9)$, so that leg moves the readout \emph{and} the placement --- and the resulting
attribution flips sign with the choice of rung on all four DeepLabV3 conditions. The
$M{=}2$ rung must be $r_2@0.1$ or the decomposition means nothing. We note this because
the error is exactly the one this subsection accuses the $5.6$--$17.6\%$ figure of: a
comparison that moves more than one thing and is then read as though it moved one. A
decomposition must move one thing at a time --- including the decomposition that accuses
someone else of not doing so.

\paragraph{The dense estimator recovers a cell the band loses.} Against SDC, the
dense rule wins $\mathbf{7/8}$ under \emph{both} risks, where the $M=2$ band wins
$6/8$: it takes back \texttt{clipseg-target}/VOC on overlap ($0.161$ vs SDC's
$0.163$; paired bootstrap $0.84$) and on the penalized boundary risk ($86.7$ vs
$87.5$; bootstrap $0.69$ --- directional, not significant). The one cell that
survives all of it is \texttt{clipseg-general}/VOC, the condition where \emph{both}
pathologies of \S\ref{sec:results} are worst ($67\%$ detection failures, $7.0\%$
nesting leak).

\begin{remark}[The degenerate tail is an asset, not the liability we called it]
\S\ref{sec:threshold} argues that the untrimmed $\int_0^1 \HDF(Y_t,\Yhat)\,dt$ is
dominated by degenerate level sets and that the estimand must be trimmed. The dense
rule's nodes $(m-\tfrac12)/32$ span $(0.016, 0.984)$ --- essentially untrimmed, tails
included --- and it \emph{wins}. Remark~\ref{rem:atom} identifies what the tail
contributes: an explicit $(1-\max_i p_i)\cdot\diag$ term, i.e.\ a max-softmax signal.
We framed that as a contaminant. It is better read as a \emph{second, free channel of
information}, fused into the score by the integral itself: measured
$\rho(\text{dense},\ \text{aMSP}) = 0.54$--$0.75$, correlated but far from degenerate.
The trimming argument should be weakened accordingly: trimming fixes the $M=2$
\emph{estimator}'s saturation, but it is not a property the \emph{estimand} needs.
\end{remark}

\paragraph{Cost.} $M=32$ costs $32$ distance transforms per present class against the
band's $2$ --- a $16\times$ constant, still $O(MCN)$ and still milliseconds per image,
with no retraining and no second forward pass. For a multi-class deployment the
trade is clearly worth taking; for a binary one it is not.

\subsection{Open items --- experiments that must be run}\label{sec:openitems}

We list these because the paper's conclusions are conditional on them.

\begin{enumerate}
\item \textbf{The penalized boundary risk --- \emph{run, and it degraded the
headline}.} Both conventions are now reported (Tables~\ref{tab:boundary}
and~\ref{tab:penalized}). The band's $8/8$ win over SDC under
$\HDF^{\mathrm{drop}}$ falls to $6/8$ under $\HDF^{\mathrm{pen}}$. The
distance-vs-area ablation is unaffected ($8/8$ under both). This item is closed;
Findings 2--3 carry the result.

\item \textbf{The surrogate ablation ($r_2$ vs.\ band width) --- \emph{run; the bridge
holds, and the estimator the theory derives is also the better one, $24/24$}.}
Remark~\ref{rem:nobridge} shows Prop.~\ref{prop:band} cannot bridge
$\HDF(\Ylo,\Yhi)$ to $r_2$, because $\HDF$ is a percentile and not a metric. The
empirical bridge nevertheless holds: the rank correlation between the deployed band
width and the $r_2$ it stands in for is
\begin{equation}
\rho(\text{band},\, r_2) \in [0.974,\ 0.992] \quad\text{across all eight conditions.}
\end{equation}
So the band width is, in practice, the two-point quadrature estimator the theory
derives --- established by measurement, not by the broken inequality. This is the
honest form of the claim, and it closes the gap Remark~\ref{rem:nobridge} opens.

\emph{The bridge, however, is not the interesting half of this ablation. The head-to-head
is.} At \emph{matched} nodes $(\alpha, 1-\alpha) = (0.1, 0.9)$, \textbf{$r_2$
\eqref{eq:r2} beats the deployed band width \eqref{eq:bandsym} on every condition under
every risk: $8/8$ under overlap, $8/8$ under $\HDF^{\mathrm{drop}}$, $8/8$ under
$\HDF^{\mathrm{pen}}$ --- $24/24$.} This is as controlled as the distance-vs-area
ablation and for the same reason: same $\alpha$, same two thresholds, same pooled
bidirectional $\HDF$, same present-class mean, same $-\diag$ floors. \emph{The only thing
that differs is the functional} --- $\HDF(\Ylo,\Yhi)$ against
$\tfrac12[\HDF(\Ylo,\Yhat) + \HDF(\Yhi,\Yhat)]$ --- so the ranking difference is
attributable to the estimator alone. It is free: no extra node, no extra distance
transform, no new tuned parameter.

\emph{The honest bound on it.} $24/24$ is a count of \emph{signs}, not of significances,
and the three risks are not three independent tests --- they share the same images.
Against a $10$k paired image bootstrap under the penalized boundary risk, the margin's CI
excludes zero on the \emph{four} CLIPSeg conditions ($+7.07$, $+9.23$, $+3.23$,
$+2.32\%$ of band AURC) and not on the four DeepLabV3 ones ($+0.27$, $+1.09$, $+1.93$,
$+2.23\%$), which are positive but null. So: $8/8$ in sign, $4/8$ significant.

This also closes the paper's one theory--practice gap, and more cleanly than the measured
$\rho$ does. Remark~\ref{rem:nobridge} exists because $\HDF$ is a percentile and no
inequality connects the band to the $r_2$ it stands in for. \emph{Deploy $r_2$ and there
is nothing to bridge} --- and we now know that deploying it costs nothing and ranks
strictly better. The band width earned its place by being cheap to compute and easy to
draw; on this evidence it has no other claim.

\emph{The same measurement condemns the rank-anchored ablation
\eqref{eq:rankanchor}}: against the dense $M=32$ reference, the rank-anchored band
scores \emph{below} the constant-threshold band in $6/8$ conditions and the
rank-anchored $r_2$ below $r_2$ in $8/8$. It is a worse estimator of the level-set
risk on real data, not merely on the synthetic populations of
\S\ref{sec:threshold}.

\item \textbf{The quadrature ablation ($M$) --- \emph{run, audited, recomputed,
unchanged}.} See \S\ref{sec:quadrature} and Remark~\ref{rem:mladdercontam}. The ladder
was measured; an audit then found an aggregation defect that dropped a hallucinated class
instead of saturating it, and identified a mechanism by which it could have fired on
exactly the dense rules and not the cheap ones, making the ladder a comparison of rules
over different sets of classes on the same image. Every number was withheld. \emph{The
recomputation has landed and the ladder reproduces to the digit.} The defect moved $52$
scores in total --- all on the two genuine $21$-way softmax conditions, all at
$\alpha = 0.3$, none in the ladder and none at the deployed $\alpha = 0.1$. The results
therefore stand as reported: denser quadrature is decisively better on the multi-class
benchmark and flat on the binary one, and the dense estimator recovers
\texttt{clipseg-target}/VOC from SDC, taking the count to $7/8$ where the $M=2$ band
manages $6/8$. \emph{What the recomputation did not do is rescue our explanation of the
$5.6$--$17.6\%$}: \S\ref{sec:quadrature} decomposes it and finds the node-count leg
indistinguishable from zero at the top of the range. This item is \textbf{closed}. The
explanation is not, and item 4 is where it died.

\item \textbf{Vertex-aware nodes on multi-class (\S\ref{sec:derivednodes}) --- \emph{run.
They lose. The fourth adaptive node rule in this paper to fail.}}
Observation~\ref{obs:vshape} places the vertex near $m_c$, the smallest probability
inside the prediction --- measured at $m_c = 0.055$ for a $21$-way softmax (against
$t^\star = 0.060$), and at $\tfrac12$ for binary. \eqref{eq:derivednodes} then prescribes
nodes $\approx (0.028, 0.528)$ on VOC --- nowhere near the deployed $(0.1, 0.9)$ or the
binary-derived $(0.25, 0.75)$, both of which put \emph{both} nodes on the same side of
the vertex and never straddle it. We billed this as the open question that mattered most,
and predicted that if Remark~\ref{rem:straddle} were right, vertex-aware $M=2$ should
recover most of what dense $M=32$ buys on VOC at $1/16$ of the cost.

\emph{It has now been run on the real conditions, and the prediction fails.} Against the
deployed band under the penalized boundary risk:

\begin{center}
\small
\begin{tabular}{lc}
\toprule
$M=2$ rule & beats the deployed band \\
\midrule
dense $M{=}32$ (reference)                          & $\mathbf{8/8}$ \\
\texttt{mid2}, nodes $(0.25, 0.75)$                 & $6/8$ \\
\texttt{vtx2}, vertex-aware \eqref{eq:derivednodes} & $\mathbf{4/8}$ --- \emph{worst of the three} \\
\bottomrule
\end{tabular}
\end{center}

\texttt{mid2} --- the fixed, unadaptive rule \texttt{vtx2} was designed to improve on ---
beats it $6/8$ head-to-head. And on the only two conditions where the vertex-aware rule
has anything to repair (\texttt{deeplabv3-external}/VOC and \texttt{deeplabv3-target}/VOC,
the only two whose band fails to straddle the vertex; see Remark~\ref{rem:straddle}),
it is \emph{worse than
the band it was meant to repair}: $49.35$ against $47.88$, and $48.46$ against $46.38$.
It fails hardest exactly where it was supposed to win.

This is the \textbf{fourth} principled-looking adaptive node rule here to be derived,
look compelling, and fail on measurement --- after rank-anchoring \eqref{eq:rankanchor},
histogram importance sampling, and the $\sqrt{V_L/V_R}$ allocation. Three of the four
have the common mechanism \S\ref{sec:threshold} identifies: they move nodes \emph{toward}
where the probability values sit, and the integral's mass is in the tails.
\texttt{vtx2} is at least consistent with it --- it puts $94.5\%$ of the weight on a node
at $\approx 0.53$ and abandons the right tail --- but we have not isolated that mechanism
here and do not claim it. What we claim is the count: four for four. \emph{That $m_c$ is
an estimate of the vertex and not the vertex (Remark~\ref{rem:vshapefails}) is no longer
the binding objection; the rule was tested, and lost.}

The consequence for the paper is the \emph{opposite} of the warning we issued here.
Earlier drafts said: ``until this is run, our $M=2$ results should be read as using
demonstrably mis-placed nodes on the multi-class benchmark.'' \textbf{That warning is
withdrawn.} On real multi-class data $(0.1, 0.9)$ beats the nodes our own integrand
analysis prescribes. The rule is also measured against the exact integral in
\texttt{scripts/vertex\_nodes.py} --- which, until an audit caught it, measured the
\emph{retracted} floor rule ($\phi = 1/C$) rather than \eqref{eq:derivednodes}'s
$\hat t = m_c$, and now measures both so the difference stays visible. This item is
\textbf{closed, negatively}.

\item \textbf{Threshold selection (\S\ref{sec:threshold}) --- half closed, and the
half that matters is still open.} \emph{The grid extension this item demanded has
been run.} The evaluation now sweeps
$\alpha\in\{0.01,0.05,0.1,0.15,0.2,0.3\}$ on all eight conditions, and the sweep is
\textbf{bracketed}: mean gap-closure is no longer rising at the grid edge but turns
over at $\alpha=0.2$ under all three risks ($0.723$ at $0.2$ against $0.719$ at
$0.3$ under overlap; $0.683$ against $0.680$ under $\HDF^{\mathrm{drop}}$; $0.692$
against $0.684$ under $\HDF^{\mathrm{pen}}$). The band-beats-SDC counts across the
grid:

\begin{center}
\small
\begin{tabular}{lcccccc}
\toprule
risk & $0.01$ & $0.05$ & $0.1$ & $0.15$ & $0.2$ & $0.3$ \\
\midrule
overlap                & $4/8$ & $5/8$ & $6/8$ & $6/8$ & $6/8$ & $6/8$ \\
$\HDF^{\mathrm{drop}}$ & $5/8$ & $6/8$ & $\mathbf{8/8}$ & $\mathbf{8/8}$ & $\mathbf{8/8}$ & $7/8$ \\
$\HDF^{\mathrm{pen}}$  & $4/8$ & $6/8$ & $6/8$ & $6/8$ & $6/8$ & $6/8$ \\
\bottomrule
\end{tabular}
\end{center}

\textbf{The core admission is unchanged and we restate it rather than retire it:
$\alpha=0.1$ was chosen on the test set, and the headline is \emph{not} robust to
it.} At $\alpha=0.05$ the band-beats-SDC count under $\HDF^{\mathrm{drop}}$ falls
from $8/8$ to $6/8$ --- a one-step move on the grid costs the headline two cells.
Bracketing the optimum does not make the selection honest; it only means we now know
that the value we picked is not even the sweep's argmax, which is $\alpha=0.2$ on
mean gap-closure under every risk. Nor is the enclosure uniform: per condition the
argmax still lands on a grid edge in four of eight cells ($\alpha=0.3$ three times,
$\alpha=0.01$ once, under each risk). What remains required is the part that was
always the point: select $\alpha$ on a validation split carved from the
\emph{training} data, not from the test set. Until that is done, treat $\alpha$ as
tuned.
\end{enumerate}

\section{Related work}

\paragraph{Selective prediction.} Thresholding a fixed network's confidence remains
the dominant baseline \cite{selectivenet}; well-tuned thresholds often match
training-time methods \cite{feng}; MSP is strong \cite{jaeger}. Calibration is not
abstention \cite{zhu}, which is why we report AURC and keep ECE separate.

\paragraph{Uncertainty and aggregation in segmentation.} ValUES \cite{values} names
score aggregation as a neglected component; that gap has since been addressed by
Guarino et al.\ \cite{guarino}, who propose spatially aware aggregators and
benchmark them for failure detection, and by Zenk et al.\ \cite{zenk}, whose
headline finding is likewise about confidence aggregation. We state our position
precisely against them: every such aggregator operates in \emph{probability units}
on a per-pixel uncertainty map; ours is derived from the deployment metric and
returns a value in \emph{pixels}. Learned failure predictors --- ConfidNet
\cite{confidnet}, ObsNet \cite{obsnet}, FSNet \cite{fsnet} --- need internal features
and extra training, outside the post-hoc single-pass regime.

\paragraph{Metric-aligned confidence.} SDC \cite{sdc} is the direct predecessor;
\S\ref{sec:theory} places it inside our framework and delimits it.

\paragraph{Random sets and uncertain contours.} The level-set posterior is the
classical coverage-function random set \cite{goodman,molchanov}; the insight that
independent per-pixel sampling destroys contour geometry is established in the
uncertain-isocontour literature \cite{poethkow,bolin}. We import these rather than
claim them (Remark~\ref{rem:classical}).

\paragraph{Distance transforms and nested masks.} Distance transforms of probability
maps are used at \emph{training} time to reduce Hausdorff distance \cite{karimi}
(which gives three such estimators), at inference to build conformal inner/outer sets
\cite{davenport,mossina}, and to evaluate uncertainty maps \cite{buc}. What is new
here is the distance transform \emph{between two level sets of the same probability
map}, as a scalar, ground-truth-free, image-level confidence. The nested inner/outer
band is itself a first-class object in the conformal literature
\cite{crc,davenport,mossina}: that line \emph{calibrates} the band to guarantee a
property of the set; we \emph{measure its width} and use it to rank images for
abstention. Same object, different use --- and \S\ref{sec:threshold} argues the two
should be combined.

\section{Limitations}\label{sec:limitations}

\paragraph{It measures blur, not error.} The band is wide when the model is unsure
and narrow when it is sure. A confidently wrong prediction --- a sharp mask in the
wrong place --- gets a narrow band and high confidence. This is shared by all
post-hoc probability-map scores and is why term (b) of \eqref{eq:decomp} cannot be
driven to zero by computation on $p$ alone.

\paragraph{The posterior is a modelling choice.} Nested level sets model ambiguity as
contour dilation and erosion. Real ambiguity also includes topological disagreement,
which the comonotone coupling does not generate.

\paragraph{The threshold is unresolved.} \S\ref{sec:threshold} is a diagnosis, not
yet a validated fix.

\paragraph{Scope.} Natural-image benchmarks only, two architectures, no medical data
and no OOD split --- while the motivation is medical and human-in-the-loop, and SDC
was validated on six medical tasks with OOD. This is the largest gap between our
claims and our evidence.

\begin{thebibliography}{99}
\small
\bibitem{sdc} Borges, B. L. C., Pacheco, B. M., \& Silva, D. \emph{Soft Dice
Confidence: A Near-Optimal Confidence Estimator for Selective Prediction in Semantic
Segmentation}. arXiv:2402.10665.
\bibitem{sam} Kirillov, A., et al. \emph{Segment Anything}. ICCV 2023.
\bibitem{values} Kahl, K., et al. \emph{ValUES: A Framework for Systematic Validation
of Uncertainty Estimation in Semantic Segmentation}. ICLR 2024 (oral).
arXiv:2401.08501.
\bibitem{guarino} Guarino, A., Winklmayr, L., Franzen, D., et al. \emph{Better than
Average: Spatially-Aware Aggregation of Segmentation Uncertainty Improves Downstream
Performance}. CVPR 2026.
\bibitem{crc} Angelopoulos, A. N., et al. \emph{Conformal Risk Control}. ICLR 2024.
arXiv:2208.02814.
\bibitem{zenk} Zenk, M., et al. \emph{Comparative Benchmarking of Failure Detection
Methods in Medical Image Segmentation: Unveiling the Role of Confidence Aggregation}.
Medical Image Analysis 2025. arXiv:2406.03323.
\bibitem{zhu} Zhu, F., et al. \emph{Rethinking Confidence Calibration for Failure
Prediction}. ECCV 2022.
\bibitem{selectivenet} Geifman, Y., \& El-Yaniv, R. \emph{SelectiveNet}. ICML 2019.
\bibitem{feng} Feng, L., Ahmed, M. O., Hajimirsadeghi, H., \& Abdi, A. \emph{Towards
Better Selective Classification}. ICLR 2023. arXiv:2206.09034.
\bibitem{jaeger} Jaeger, P. F., et al. \emph{A Call to Reflect on Evaluation
Practices for Failure Detection in Image Classification}. ICLR 2023.
\bibitem{confidnet} Corbi\`ere, C., et al. \emph{Addressing Failure Prediction by
Learning Model Confidence}. NeurIPS 2019.
\bibitem{obsnet} Besnier, V., et al. \emph{Triggering Failures (ObsNet)}. ICCV 2021.
\bibitem{fsnet} Rahman, Q. M., et al. \emph{FSNet}. IEEE RA-L 2022.
\bibitem{karimi} Karimi, D., \& Salcudean, S. E. \emph{Reducing the Hausdorff
Distance in Medical Image Segmentation with Convolutional Neural Networks}. IEEE TMI
39(2):499--513, 2020.
\bibitem{gen} Liu, X., et al. \emph{GEN: Pushing the Limits of Softmax-Based OOD
Detection}. CVPR 2023.
\bibitem{doctor} Granese, F., et al. \emph{DOCTOR: A Simple Method for Detecting
Misclassification Errors}. NeurIPS 2021.
\bibitem{metaseg} Rottmann, M., et al. \emph{Prediction Error Meta Classification in
Semantic Segmentation (MetaSeg)}. IJCNN 2020.
\bibitem{holder} Holder, C. J., \& Shafique, M. \emph{Efficient Uncertainty
Estimation in Semantic Segmentation via Distillation}. ICCVW 2021. (MMMC.)
\bibitem{buc} Zeevi, T., et al. \emph{Spatially-Aware Evaluation of Segmentation
Uncertainty}. UNCV @ CVPR 2025.
\bibitem{goodman} Goodman, I. R. \emph{Fuzzy sets as equivalence classes of random
sets}. 1982.
\bibitem{molchanov} Molchanov, I. \emph{Theory of Random Sets}. Springer, 2nd ed.,
2017.
\bibitem{poethkow} P\"othkow, K., Weber, B., \& Hege, H.-C. \emph{Probabilistic
Marching Cubes}. EuroVis 2011.
\bibitem{bolin} Bolin, D., \& Lindgren, F. \emph{Excursion and contour uncertainty
regions for latent Gaussian models}. JRSS-B 77(1):85--106, 2015.
\bibitem{davenport} Davenport, S., et al. \emph{Conformal prediction sets for
segmentation}. arXiv:2410.03406.
\bibitem{mossina} Mossina, L., \& Friedrich, C. \emph{Conformal morphological
prediction sets for segmentation}. MICCAI 2025.
\end{thebibliography}

\end{document}
